\documentclass[11pt,a4paper]{article}
\usepackage[utf8]{inputenc}
\usepackage[margin=0.75in]{geometry}
\usepackage{amsmath,amssymb,amsfonts,amsthm}
\usepackage{enumitem}
\usepackage{booktabs}
\usepackage{tabularx}
\usepackage{xcolor}
\usepackage{titlesec}
\usepackage{tcolorbox}
\usepackage[colorlinks=true,linkcolor=blue,urlcolor=blue,citecolor=blue]{hyperref}

% Define Palette
\definecolor{primary}{RGB}{20, 45, 85}
\definecolor{secondary}{RGB}{41, 128, 185}
\definecolor{badgebg}{RGB}{240, 243, 246}
\definecolor{badgetext}{RGB}{40, 40, 40}

% Section Styling
\titleformat{\section}{\Large\bfseries\color{primary}}{\thesection}{1em}{}[\titlerule]
\titleformat{\subsection}{\large\bfseries\color{secondary}}{\thesubsection}{1em}{}

% Custom Question Badge Environment
\newcommand{\pyqheader}[4]{%
    \noindent\begin{tcolorbox}[
        colback=badgebg,
        colframe=secondary,
        arc=2mm,
        boxrule=0.6pt,
        top=1.5mm,bottom=1.5mm,left=2mm,right=2mm
    ]
        \textbf{\textcolor{primary}{JAM #1}} \hfill \textbf{Q.#2} \hfill \textsc{\textbf{#3}} \hfill \textbf{[#4 Mark\ifnum#4>1s\fi]}
    \end{tcolorbox}\vspace{-1mm}
}

\begin{document}

\begin{center}
    {\huge \textbf{\color{primary}IIT JAM Mathematics (MA)}}\\[2mm]
    {\Large \textbf{Topic-Wise Question Bank (2022 -- 2026)}}\\[1.5mm]
    {\normalsize Categorized by Official Syllabus with Complete Master Answer Key}
\end{center}

\tableofcontents
\vspace{5mm}
\hrule
\vspace{5mm}

% ==============================================================================
% SECTION 1: REAL ANALYSIS
% ==============================================================================
\section{Real Analysis}

\subsection{Sequences and Series of Real Numbers}

\pyqheader{2026}{1}{MCQ}{1}
For each positive integer $n$, let $x_{n}=1-(-1)^{n}+\frac{1}{n}$ and $y_{n}=\left(1+\frac{1}{2n}\right)^{3n}$. Which ONE of the following statements about the sequences $(x_{n})$ and $(y_{n})$ is TRUE?
\begin{enumerate}[label=(\Alph*)]
    \item $(x_{n})$ and $(y_{n})$ are convergent.
    \item $(x_{n})$ is not convergent and $(y_{n})$ is not convergent.
    \item $(x_{n})$ is convergent and $(y_{n})$ is not convergent.
    \item $(x_{n})$ is not convergent and $(y_{n})$ is convergent.
\end{enumerate}

\pyqheader{2026}{2}{MCQ}{1}
For each positive integer $n$, define $x_{n}=(-1)^{n}$. Which ONE of the following statements about the sequence $(x_{n})$ is FALSE?
\begin{enumerate}[label=(\Alph*)]
    \item There exists $\epsilon>0$ such that $|x_{n}-1| < \epsilon$ for all positive integers $n$.
    \item There exists $\epsilon>0$ such that for all $M>0$ there exists a positive integer $N>M$ for which $|x_{N}-1|>\epsilon$.
    \item For all $\epsilon>0$ and $M>0$ there exists a positive integer $N$ such that $N>M$ and $|x_{N}-1|<\epsilon$.
    \item For all $\epsilon>0$ and $M>0$ there exists a positive integer $N$ such that $N>M$ and $|x_{N}-1|>\epsilon$.
\end{enumerate}

\pyqheader{2026}{11}{MCQ}{2}
Given a sequence of real numbers $(a_{n})$, define $b_{n}=\begin{cases}a_{n},& \text{if } a_{n}\ge 0\\ 0,& \text{if } a_{n}<0\end{cases}$ and $c_{n}=\begin{cases}a_{n},& \text{if } a_{n}< 0\\ 0,& \text{if } a_{n}\ge 0\end{cases}$ for each positive integer $n$. Which ONE of the following statements is TRUE?
\begin{enumerate}[label=(\Alph*)]
    \item If $\sum_{n=1}^{\infty}a_{n}$ does not converge, then $\sum_{n=1}^{\infty}b_{n}$ does not converge.
    \item If $\sum_{n=1}^{\infty}a_{n}$ converges but $\sum_{n=1}^{\infty}|a_{n}|$ does not converge, then $\sum_{n=1}^{\infty}b_{n}$ converges.
    \item If $\sum_{n=1}^{\infty}a_{n}$ converges but $\sum_{n=1}^{\infty}|a_{n}|$ does not converge, then $\sum_{n=1}^{\infty}b_{n}$ does not converge.
    \item If $\sum_{n=1}^{\infty}a_{n}$ converges, then $\sum_{n=1}^{\infty}b_{n}$ converges or $\sum_{n=1}^{\infty}c_{n}$ converges.
\end{enumerate}

\pyqheader{2026}{12}{MCQ}{2}
For which one of the following pairs of values of $a$ and $b$, does the series $\sum_{k=1}^{\infty}\frac{(3x)^{k}}{2^{\sqrt{k}}(1-x)^{k}}$ converge for all $x\in(a,b)$?
\begin{enumerate}[label=(\Alph*)]
    \item $a=-\frac{1}{4}$ and $b=\frac{1}{2}$
    \item $a=-1$ and $b=\frac{1}{5}$
    \item $a=-\frac{1}{2}$ and $b=\frac{1}{4}$
    \item $a=-\frac{1}{4}$ and $b=\frac{1}{3}$
\end{enumerate}

\pyqheader{2026}{29}{MCQ}{2}
For each positive integer $n$, let $x_{n}=1+\frac{1}{2}+\cdots+\frac{1}{n}-\log_{e}n$ and $y_{n}=\int_{1}^{n}\frac{\cos t}{t^{2}}dt$. Which ONE of the following statements about $(x_{n})$ and $(y_{n})$ is TRUE?
\begin{enumerate}[label=(\Alph*)]
    \item $(x_{n})$ and $(y_{n})$ are convergent.
    \item $(x_{n})$ is convergent and $(y_{n})$ is not convergent.
    \item $(x_{n})$ is not convergent and $(y_{n})$ is convergent.
    \item $(x_{n})$ is not convergent and $(y_{n})$ is not convergent.
\end{enumerate}

\pyqheader{2026}{32}{MSQ}{2}
Which of the following statements is/are TRUE?
\begin{enumerate}[label=(\Alph*)]
    \item There exists a monotone sequence that does not converge but has a convergent subsequence.
    \item There exists a sequence that has a bounded subsequence but does not have any convergent subsequence.
    \item There exists a sequence $(x_{n})$ such that given any positive integer $m$, $(x_{n})$ has a subsequence converging to $m$.
    \item There exists a sequence $(x_{n})$ such that $(|x_{n+1}-x_{n}|)$ converges to 0 but $(x_{n})$ does not converge.
\end{enumerate}

\pyqheader{2026}{41}{NAT}{1}
The radius of convergence of the series $\sum_{n=1}^{\infty}\frac{(n!)^{2}}{(2n)!}(\log_{e}n)^{-1}x^{n}$ is \underline{\hspace{1.5cm}} (rounded off to one decimal place).

\pyqheader{2026}{42}{NAT}{1}
$\int_{0}^{1}\left(\sum_{k=1}^{\infty}\frac{(\log_{e}2)^{k}x^{k-1}}{(k-1)!}\right)dx=$ \underline{\hspace{1.5cm}} (rounded off to one decimal place).

\pyqheader{2026}{43}{NAT}{1}
$\lim_{n\rightarrow\infty}\left(\frac{n^{2}+1}{\sqrt{n^{6}+1}}+\cdots+\frac{n^{2}+n}{\sqrt{n^{6}+n}}\right)=$ \underline{\hspace{1.5cm}} (rounded off to one decimal place).

\pyqheader{2025}{1}{MCQ}{1}
The sum of the infinite series $\sum_{n=1}^{\infty}(-1)^{n+1}\frac{\pi^{2n+1}}{2^{2n+1}(2n)!}$ is equal to:
\begin{enumerate}[label=(\Alph*)]
    \item $-\pi$
    \item $\frac{\pi}{4}$
    \item $\frac{\pi}{2}$
    \item $-\frac{\pi}{4}$
\end{enumerate}

\pyqheader{2025}{18}{MCQ}{2}
For $n\in\mathbb{N}$, define $x_{n}=(-1)^{n}\cos\frac{1}{n}$ and $y_{n}=\sum_{k=1}^{n}\frac{1}{n+k}$. Then, which one of the following is TRUE?
\begin{enumerate}[label=(\Alph*)]
    \item $\sum_{n=1}^{\infty} x_n$ converges, and $\sum_{n=1}^{\infty} y_n$ does NOT converge
    \item $\sum_{n=1}^{\infty} x_n$ does NOT converge, and $\sum_{n=1}^{\infty} y_n$ converges
    \item $\sum_{n=1}^{\infty} x_n$ converges, and $\sum_{n=1}^{\infty} y_n$ converges
    \item $\sum_{n=1}^{\infty} x_n$ does NOT converge, and $\sum_{n=1}^{\infty} y_n$ does NOT converge
\end{enumerate}

\pyqheader{2025}{19}{MCQ}{2}
Let $x_{1}=\frac{5}{2}$. For $n\in\mathbb{N}$, define $x_{n+1}=\frac{1}{5}(x_{n}^{2}+6)$. Then, which one of the following is TRUE?
\begin{enumerate}[label=(\Alph*)]
    \item $(x_{n})$ is an increasing sequence, and $(x_{n})$ is NOT a bounded sequence
    \item $(x_{n})$ is NOT an increasing sequence, and $(x_{n})$ is NOT a bounded sequence
    \item $(x_{n})$ is NOT a decreasing sequence, and $(x_{n})$ is a bounded sequence
    \item $(x_{n})$ is a decreasing sequence, and $(x_{n})$ is a bounded sequence
\end{enumerate}

\pyqheader{2025}{20}{MCQ}{2}
Let $x_{1}=2$ and $x_{n+1}=2+\frac{1}{2x_{n}}$ for all $n\in\mathbb{N}$. Then, which one of the following is TRUE?
\begin{enumerate}[label=(\Alph*)]
    \item $x_{n+1}\ge\frac{4}{x_{n}}$ for all $n\in\mathbb{N}$, and $(x_{n})$ is a Cauchy sequence
    \item $x_{n+1}<\frac{4}{x_{n}}$ for some $n\in\mathbb{N}$, and $(x_{n})$ is a Cauchy sequence
    \item $x_{n+1}\ge\frac{4}{x_{n}}$ for all $n\in\mathbb{N}$, and $(x_{n})$ is NOT a Cauchy sequence
    \item $x_{n+1}<\frac{4}{x_{n}}$ for some $n\in\mathbb{N}$, and $(x_{n})$ is NOT a Cauchy sequence
\end{enumerate}

\pyqheader{2025}{21}{MCQ}{2}
For $n\in\mathbb{N}$, define $x_{n}=(-1)^{n}\frac{3^{n}}{n^{3}}$ and $y_{n}=(4^{n}+(-1)^{n}3^{n})^{1/n}$. Then, which one of the following is TRUE?
\begin{enumerate}[label=(\Alph*)]
    \item $(x_{n})$ has a convergent subsequence, and NO subsequence of $(y_{n})$ is convergent
    \item NO subsequence of $(x_{n})$ is convergent, and $(y_{n})$ has a convergent subsequence
    \item $(x_{n})$ has a convergent subsequence, and $(y_{n})$ has a convergent subsequence
    \item NO subsequence of $(x_{n})$ is convergent, and NO subsequence of $(y_{n})$ is convergent
\end{enumerate}

\pyqheader{2025}{25}{MCQ}{2}
Let $x_{1}=1$. For $n\in\mathbb{N}$, define $x_{n+1}=\left(\frac{1}{2}+\frac{\sin^{2}n}{n}\right)x_{n}$. Then, which one of the following is TRUE?
\begin{enumerate}[label=(\Alph*)]
    \item $\sum_{n=1}^{\infty}x_{n}$ converges
    \item $\sum_{n=1}^{\infty}x_{n}$ does NOT converge
    \item $\sum_{n=1}^{\infty}x_{n}^{2}$ does NOT converge
    \item $\sum_{n=1}^{\infty}x_{n}x_{n+1}$ does NOT converge
\end{enumerate}

\pyqheader{2025}{26}{MCQ}{2}
Let $x_{1}>0$. For $n\in\mathbb{N}$, define $x_{n+1}=x_{n}+4$. If $\lim_{n\rightarrow\infty}\left(\frac{1}{x_{2}x_{3}}+\frac{1}{x_{3}x_{4}}+\cdots+\frac{1}{x_{n+1}x_{n+2}}\right)=\frac{1}{24}$, then the value of $x_{1}$ is equal to:
\begin{enumerate}[label=(\Alph*)]
    \item 1
    \item 2
    \item 3
    \item 4
\end{enumerate}

\pyqheader{2025}{36}{MSQ}{2}
For $n\in\mathbb{N}$, let $x_{n}=\sum_{k=1}^{n}\frac{k}{n^{2}+k}$. Then, which of the following is/are TRUE?
\begin{enumerate}[label=(\Alph*)]
    \item The sequence $(x_{n})$ converges
    \item The series $\sum_{n=1}^{\infty}x_{n}$ converges
    \item The series $\sum_{n=1}^{\infty}x_{n}$ does NOT converge
    \item The series $\sum_{n=1}^{\infty}x_{n}^{n}$ converges
\end{enumerate}

\pyqheader{2025}{41}{NAT}{1}
The radius of convergence of the power series $\sum_{n=1}^{\infty}\frac{(x+\frac{1}{4})^{n}}{(-2)^{n}n^{2}}$ about $x=-\frac{1}{4}$ is equal to \underline{\hspace{1.5cm}} (rounded off to two decimal places).

\pyqheader{2025}{54}{NAT}{2}
$\sum_{n=1}^{\infty}n\left(\frac{3}{4}\right)^{2(n-1)}=$ \underline{\hspace{1.5cm}} (rounded off to two decimal places).

\pyqheader{2024}{14}{MCQ}{2}
For $n\in\mathbb{N}$, let $a_{n}=\frac{1}{(3n+2)(3n+4)}$ and $b_{n}=\frac{n^{3}+\cos(3^{n})}{3^{n}+n^{3}}$. Then, which one of the following is TRUE?
\begin{enumerate}[label=(\Alph*)]
    \item $\sum_{n=1}^{\infty} a_n$ is convergent but $\sum_{n=1}^{\infty} b_n$ is divergent
    \item $\sum_{n=1}^{\infty} a_n$ is divergent but $\sum_{n=1}^{\infty} b_n$ is convergent
    \item Both $\sum_{n=1}^{\infty} a_n$ and $\sum_{n=1}^{\infty} b_n$ are divergent
    \item Both $\sum_{n=1}^{\infty} a_n$ and $\sum_{n=1}^{\infty} b_n$ are convergent
\end{enumerate}

\pyqheader{2024}{20}{MCQ}{2}
Define the sequences $\{a_{n}\}_{n=3}^{\infty}$ and $\{b_{n}\}_{n=3}^{\infty}$ as $a_{n}=(\log n+\log\log n)^{\log n}$ and $b_{n}=n^{(1+\frac{1}{\log n})}$. Which one of the following is TRUE?
\begin{enumerate}[label=(\Alph*)]
    \item $\sum_{n=3}^{\infty} \frac{1}{a_n}$ is convergent but $\sum_{n=3}^{\infty} \frac{1}{b_n}$ is divergent
    \item $\sum_{n=3}^{\infty} \frac{1}{a_n}$ is divergent but $\sum_{n=3}^{\infty} \frac{1}{b_n}$ is convergent
    \item Both $\sum_{n=3}^{\infty} \frac{1}{a_n}$ and $\sum_{n=3}^{\infty} \frac{1}{b_n}$ are divergent
    \item Both $\sum_{n=3}^{\infty} \frac{1}{a_n}$ and $\sum_{n=3}^{\infty} \frac{1}{b_n}$ are convergent
\end{enumerate}

\pyqheader{2024}{21}{MCQ}{2}
For $p,q,r\in\mathbb{R}$, $r\ne 0$ and $n\in\mathbb{N}$, let $a_{n}=p^{n}n^{q}\left(\frac{n}{n+2}\right)^{n^{2}}$ and $b_{n}=\frac{n^{n}}{n!r^{n}}\left(\sqrt{\frac{n+2}{n}}\right)$. Then:
\begin{enumerate}[label=(\Alph*)]
    \item If $1<p<e^{2}$ and $q>1$, then $\sum_{n=1}^{\infty} a_{n}$ is convergent
    \item If $e^{2}<p<e^{4}$ and $q>1$, then $\sum_{n=1}^{\infty} a_{n}$ is convergent
    \item If $1<r<e$, then $\sum_{n=1}^{\infty} b_{n}$ is convergent
    \item If $\frac{1}{e}<r<1$, then $\sum_{n=1}^{\infty} b_{n}$ is convergent
\end{enumerate}

\pyqheader{2024}{28}{MCQ}{2}
For $n\ge 3$, let a regular $n$-sided polygon $P_{n}$ be circumscribed by a circle of radius $R_{n}$ and let $r_{n}$ be the radius of the circle inscribed in $P_{n}$. Then $\lim_{n\rightarrow\infty}\left(\frac{R_{n}}{r_{n}}\right)^{n^{2}}$ equals:
\begin{enumerate}[label=(\Alph*)]
    \item $e^{(\pi^{2})}$
    \item $e^{(\frac{\pi^{2}}{2})}$
    \item $e^{(\frac{\pi^{2}}{3})}$
    \item $e^{(2\pi^{2})}$
\end{enumerate}

\pyqheader{2024}{31}{MSQ}{2}
Let $\{a_{n}\}_{n=1}^{\infty}$ be a sequence of real numbers. Then, which of the following statements is/are always TRUE?
\begin{enumerate}[label=(\Alph*)]
    \item If $\sum_{n=1}^{\infty} a_{n}$ converges absolutely, then $\sum_{n=1}^{\infty} a_{n}^{2}$ converges absolutely
    \item If $\sum_{n=1}^{\infty} a_{n}$ converges absolutely, then $\sum_{n=1}^{\infty} a_{n}^{3}$ converges absolutely
    \item If $\sum_{n=1}^{\infty} a_{n}$ converges, then $\sum_{n=1}^{\infty} a_{n}^{2}$ converges
    \item If $\sum_{n=1}^{\infty} a_{n}$ converges, then $\sum_{n=1}^{\infty} a_{n}^{3}$ converges
\end{enumerate}

\pyqheader{2024}{32}{MSQ}{2}
Which of the following statements is/are TRUE?
\begin{enumerate}[label=(\Alph*)]
    \item $\sum_{n=1}^{\infty}n\log(1+\frac{1}{n^{3}})$ is convergent
    \item $\sum_{n=1}^{\infty}(1-\cos(\frac{1}{n}))\log n$ is convergent
    \item $\sum_{n=1}^{\infty}n^{2}\log(1+\frac{1}{n^{3}})$ is convergent
    \item $\sum_{n=1}^{\infty}(1-\cos(\frac{1}{\sqrt{n}}))\log n$ is convergent
\end{enumerate}

\pyqheader{2024}{39}{MSQ}{2}
For $0<\alpha<4$, define the sequence $\{x_{n}\}_{n=1}^{\infty}$ as $x_{1}=\alpha$ and $x_{n+1}+2=-x_{n}(x_{n}-4)$ for $n\in\mathbb{N}$. Which of the following is/are TRUE?
\begin{enumerate}[label=(\Alph*)]
    \item $\{x_{n}\}$ converges for at least three distinct values of $\alpha\in(0,1)$
    \item $\{x_{n}\}$ converges for at least three distinct values of $\alpha\in(1,2)$
    \item $\{x_{n}\}$ converges for at least three distinct values of $\alpha\in(2,3)$
    \item $\{x_{n}\}$ converges for at least three distinct values of $\alpha\in(3,4)$
\end{enumerate}

\pyqheader{2024}{47}{NAT}{1}
Let $a_1=1, b_1=2, c_1=3$. Define $a_{n+1}=\frac{a_n+b_n}{2}, b_{n+1}=\frac{b_n+c_n}{2}, c_{n+1}=\frac{c_n+a_n}{2}$. Then,
\[
\sum_{n=1}^{\infty}b_n c_n(a_{n+1}-a_n) + \sum_{n=1}^{\infty}(b_{n+1}c_{n+1}-b_n c_n)a_{n+1} = \underline{\hspace{1.5cm}}
\]
(rounded off to two decimal places).

\pyqheader{2024}{55}{NAT}{2}
If $a_{n}=\frac{1}{n^{3}+1}+\frac{2^{2}}{n^{3}+2}+\cdots+\frac{n^{2}}{n^{3}+n}$, then the sequence $\{a_{n}\}$ converges to \underline{\hspace{1.5cm}} (rounded off to two decimal places).

\pyqheader{2024}{59}{NAT}{2}
Let $a_{n}=\frac{1}{n^{n-1}}\sum_{k=0}^{n}\frac{n!}{k!(n-k)!}\frac{n^{k}}{k+1}$ and $\beta=\lim_{n\rightarrow\infty}a_{n}$. Then, the value of $\log\beta$ equals \underline{\hspace{1.5cm}} (rounded off to two decimal places).

\pyqheader{2023}{9}{MCQ}{1}
Let $a_{n}=(1+\frac{1}{n})^{n}$ and $b_{n}=n\cos(\frac{n!\pi}{2^{10}})$ for $n\in\mathbb{N}$. Then:
\begin{enumerate}[label=(\Alph*)]
    \item $(a_{n})$ is convergent and $(b_{n})$ is bounded
    \item $(a_{n})$ is NOT convergent and $(b_{n})$ is bounded
    \item $(a_{n})$ is convergent and $(b_{n})$ is unbounded
    \item $(a_{n})$ is NOT convergent and $(b_{n})$ is unbounded
\end{enumerate}

\pyqheader{2023}{10}{MCQ}{1}
Let $(a_{n})$ be a sequence of real numbers defined by $a_{n}=\begin{cases}1,& \text{if } n \text{ is prime}\\ -1,& \text{if } n \text{ is not prime}\end{cases}$. Let $b_{n}=\frac{a_{n}}{n}$ for $n\in\mathbb{N}$. Then:
\begin{enumerate}[label=(\Alph*)]
    \item both $(a_{n})$ and $(b_{n})$ are convergent
    \item $(a_{n})$ is convergent but $(b_{n})$ is NOT convergent
    \item $(a_{n})$ is NOT convergent but $(b_{n})$ is convergent
    \item both $(a_{n})$ and $(b_{n})$ are NOT convergent
\end{enumerate}

\pyqheader{2023}{11}{MCQ}{2}
Let $a_{n}=\sin(\frac{1}{n^{3}})$ and $b_{n}=\sin(\frac{1}{n})$ for $n\in\mathbb{N}$. Then:
\begin{enumerate}[label=(\Alph*)]
    \item both $\sum_{n=1}^{\infty} a_{n}$ and $\sum_{n=1}^{\infty} b_{n}$ are convergent
    \item $\sum_{n=1}^{\infty} a_{n}$ is convergent but $\sum_{n=1}^{\infty} b_{n}$ is NOT convergent
    \item $\sum_{n=1}^{\infty} a_{n}$ is NOT convergent but $\sum_{n=1}^{\infty} b_{n}$ is convergent
    \item both $\sum_{n=1}^{\infty} a_{n}$ and $\sum_{n=1}^{\infty} b_{n}$ are NOT convergent
\end{enumerate}

\pyqheader{2023}{19}{MCQ}{2}
Let $a_{n}=\frac{1+2^{-2}+\cdots+n^{-2}}{n}$ for $n\in\mathbb{N}$. Then:
\begin{enumerate}[label=(\Alph*)]
    \item both $(a_{n})$ and $\sum_{n=1}^{\infty} a_{n}$ are convergent
    \item $(a_{n})$ is convergent but $\sum_{n=1}^{\infty} a_{n}$ is NOT convergent
    \item both $(a_{n})$ and $\sum_{n=1}^{\infty} a_{n}$ are NOT convergent
    \item $(a_{n})$ is NOT convergent but $\sum_{n=1}^{\infty} a_{n}$ is convergent
\end{enumerate}

\pyqheader{2023}{20}{MCQ}{2}
Let $(a_{n})$ be a sequence of real numbers such that $\sum_{n=0}^{\infty}a_{n}(x-2)^{n}$ converges at $x=-5$. Then this series also converges at:
\begin{enumerate}[label=(\Alph*)]
    \item $x=9$
    \item $x=12$
    \item $x=5$
    \item $x=-6$
\end{enumerate}

\pyqheader{2023}{21}{MCQ}{2}
Let $(a_{n})$ and $(b_{n})$ be sequences of real numbers such that $|a_{n}-a_{n+1}|=\frac{1}{2^{n}}$ and $|b_{n}-b_{n+1}|=\frac{1}{\sqrt{n}}$ for $n\in\mathbb{N}$. Then:
\begin{enumerate}[label=(\Alph*)]
    \item both $(a_{n})$ and $(b_{n})$ are Cauchy sequences
    \item $(a_{n})$ is a Cauchy sequence but $(b_{n})$ need NOT be a Cauchy sequence
    \item $(a_{n})$ need NOT be a Cauchy sequence but $(b_{n})$ is a Cauchy sequence
    \item both $(a_{n})$ and $(b_{n})$ need NOT be Cauchy sequences
\end{enumerate}

\pyqheader{2023}{24}{MCQ}{2}
Suppose $f:(-1,1)\rightarrow\mathbb{R}$ is an infinitely differentiable function such that $\sum_{j=0}^{\infty}a_{j}\frac{x^{j}}{j!}$ converges to $f(x)$ for each $x\in(-1,1)$, where $a_{j}=\int_{0}^{\pi/2}\theta^{j}\cos^{j}(\tan\theta)d\theta + \int_{\pi/2}^{\pi}(\theta-\pi)^{j}\cos^{j}(\tan\theta)d\theta$ for $j\ge 0$. Then:
\begin{enumerate}[label=(\Alph*)]
    \item $f(x)=0$ for all $x\in(-1,1)$
    \item $f$ is a non-constant even function on $(-1,1)$
    \item $f$ is a non-constant odd function on $(-1,1)$
    \item $f$ is NEITHER an odd function NOR an even function on $(-1,1)$
\end{enumerate}

\pyqheader{2023}{40}{MSQ}{2}
Let $R_{1}$ and $R_{2}$ be the radii of convergence of $\sum_{n=1}^{\infty}(-1)^{n}x^{n-1}$ and $\sum_{n=1}^{\infty}(-1)^{n}\frac{x^{n+1}}{n(n+1)}$, respectively. Then:
\begin{enumerate}[label=(\Alph*)]
    \item $R_{1}=R_{2}$
    \item $R_{2}>1$
    \item $\sum_{n=1}^{\infty}(-1)^{n}x^{n-1}$ converges for all $x\in[-1,1]$
    \item $\sum_{n=1}^{\infty}(-1)^{n}\frac{x^{n+1}}{n(n+1)}$ converges for all $x\in[-1,1]$
\end{enumerate}

\pyqheader{2023}{49}{NAT}{1}
The sum of the series $\sum_{n=1}^{\infty}\frac{2n+1}{(n^{2}+1)(n^{2}+2n+2)}$ is equal to \underline{\hspace{1.5cm}} (rounded off to two decimal places).

\pyqheader{2023}{50}{NAT}{1}
$\lim_{n\rightarrow\infty}\left(1+\frac{1}{2^{n}}+\frac{1}{3^{n}}+\cdots+\frac{1}{(2023)^{n}}\right)^{\frac{1}{n}}=$ \underline{\hspace{1.5cm}} (rounded off to two decimal places).

\pyqheader{2022}{5}{MCQ}{1}
The radius of convergence of the power series $\sum_{n=1}^{\infty}\left(\frac{n^{3}}{4^{n}}\right)x^{5n}$ is:
\begin{enumerate}[label=(\Alph*)]
    \item 4
    \item $\sqrt[5]{4}$
    \item $\frac{1}{4}$
    \item $\frac{1}{\sqrt[5]{4}}$
\end{enumerate}

\pyqheader{2022}{6}{MCQ}{1}
Let $(x_{n})$ and $(y_{n})$ be defined by $x_{1}=1, y_{1}=\frac{1}{2}, x_{n+1}=\frac{x_{n}+y_{n}}{2}$, and $y_{n+1}=\sqrt{x_{n}y_{n}}$ for all $n\in\mathbb{N}$. Then:
\begin{enumerate}[label=(\Alph*)]
    \item $(x_{n})$ is convergent, but $(y_{n})$ is not convergent.
    \item $(x_{n})$ is not convergent, but $(y_{n})$ is convergent.
    \item Both $(x_{n})$ and $(y_{n})$ are convergent and $\lim_{n\rightarrow\infty}x_{n}>\lim_{n\rightarrow\infty}y_{n}$.
    \item Both $(x_{n})$ and $(y_{n})$ are convergent and $\lim_{n\rightarrow\infty}x_{n}=\lim_{n\rightarrow\infty}y_{n}$.
\end{enumerate}

\pyqheader{2022}{7}{MCQ}{1}
Suppose $a_{n}=\frac{3^{n}+3}{5^{n}-5}$ and $b_{n}=\frac{1}{(1+n^{2})^{\frac{1}{4}}}$ for $n=2,3,4,\dots$. Then:
\begin{enumerate}[label=(\Alph*)]
    \item $\sum_{n=2}^{\infty} a_{n}$ is convergent and $\sum_{n=2}^{\infty} b_{n}$ is convergent.
    \item $\sum_{n=2}^{\infty} a_{n}$ is convergent and $\sum_{n=2}^{\infty} b_{n}$ is divergent.
    \item $\sum_{n=2}^{\infty} a_{n}$ is divergent and $\sum_{n=2}^{\infty} b_{n}$ is convergent.
    \item $\sum_{n=2}^{\infty} a_{n}$ is divergent and $\sum_{n=2}^{\infty} b_{n}$ is divergent.
\end{enumerate}

\pyqheader{2022}{8}{MCQ}{1}
Consider the series $\sum_{n=1}^{\infty}\frac{1}{n^{m}(1+\frac{1}{n^{p}})}$, where $m$ and $p$ are real numbers. Under which of the following conditions does the above series converge?
\begin{enumerate}[label=(\Alph*)]
    \item $m>1$
    \item $0<m<1$ and $p>1$
    \item $0<m\le 1$ and $0\le p\le 1$
    \item $m=1$ and $p>1$
\end{enumerate}

\pyqheader{2022}{12}{MCQ}{2}
Let $T = 1-\frac{1}{2}+\frac{1}{3}-\frac{1}{4}+\dots$ and $S = 1-\frac{1}{2}-\frac{1}{4}+\frac{1}{3}-\frac{1}{6}-\frac{1}{8}+\dots$. Then:
\begin{enumerate}[label=(\Alph*)]
    \item $T=S$ and $S\ne 0$
    \item $2T=S$ and $S\ne 0$
    \item $T=2S$ and $S\ne 0$
    \item $T=S=0$
\end{enumerate}

\pyqheader{2022}{14}{MCQ}{2}
The value of $\lim_{n\rightarrow\infty}\sum_{k=2}^{n}\frac{\sqrt{n+1}-\sqrt{n}}{k(\ln k)^{2}}$ is equal to:
\begin{enumerate}[label=(\Alph*)]
    \item $\infty$
    \item 1
    \item $e$
    \item 0
\end{enumerate}

\pyqheader{2022}{28}{MCQ}{2}
For $-1\le x\le 1$, if $f(x)$ is the sum of the convergent power series $x+\frac{x^2}{2^2}+\frac{x^3}{3^2}+\dots+\frac{x^n}{n^2}+\dots$, then $f(\frac{1}{2})$ is equal to:
\begin{enumerate}[label=(\Alph*)]
    \item $\int_{0}^{1/2}\frac{\ln(1-t)}{t}dt$
    \item $-\int_{0}^{1/2}\frac{\ln(1-t)}{t}dt$
    \item $\int_{0}^{1/2}t\ln(1+t)dt$
    \item $\int_{0}^{1/2}t\ln(1-t)dt$
\end{enumerate}

\pyqheader{2022}{38}{MSQ}{2}
Let $(x_{n})$ be a sequence of real numbers. Consider $P=\{n\in\mathbb{N}: x_{n}>x_{m}\text{ for all } m>n\}$. Then:
\begin{enumerate}[label=(\Alph*)]
    \item If $P$ is finite, then $(x_{n})$ has a monotonically increasing subsequence.
    \item If $P$ is finite, then no subsequence of $(x_{n})$ is monotonically increasing.
    \item If $P$ is infinite, then $(x_{n})$ has a monotonically decreasing subsequence.
    \item If $P$ is infinite, then no subsequence of $(x_{n})$ is monotonically decreasing.
\end{enumerate}

\pyqheader{2022}{44}{NAT}{1}
The sum of the series $\sum_{n=1}^{\infty}\frac{1}{(4n-3)(4n+1)}$ is equal to \underline{\hspace{1.5cm}} (rounded off to two decimal places).

\pyqheader{2022}{55}{NAT}{2}
Let $r$ be the radius of convergence of $\frac{1}{3}+\frac{x}{5}+\frac{x^{2}}{3^{2}}+\frac{x^{3}}{5^{2}}+\frac{x^{4}}{3^{3}}+\cdots$. Then the value of $r^{2}$ is equal to \underline{\hspace{1.5cm}} (rounded off to two decimal places).

% ------------------------------------------------------------------------------
\subsection{Functions of One Real Variable}

\pyqheader{2026}{3}{MCQ}{1}
Let $f:(1,4)\rightarrow\mathbb{R}$ be a differentiable function such that $f^{\prime}(x)=(f(x)-\pi x)^{2}+\pi$ for all $x\in(1,4)$. Which ONE of the following is a possible value of $f(3)-f(2)$?
\begin{enumerate}[label=(\Alph*)]
    \item $\pi+\frac{1}{6}$
    \item $\pi-\frac{1}{6}$
    \item $\frac{\pi}{2}+\frac{1}{3}$
    \item $\frac{\pi}{2}+\frac{1}{2}$
\end{enumerate}

\pyqheader{2026}{6}{MCQ}{1}
Let $f:(1,2)\rightarrow\mathbb{R}$ be a continuous function satisfying $\lim_{x\rightarrow 1+}f(x)=3$ and $\lim_{x\rightarrow 2-}f(x)=3$. Then which ONE of the following is necessarily TRUE?
\begin{enumerate}[label=(\Alph*)]
    \item $f$ is bounded and $f$ has a maximum or a minimum but not both.
    \item $f$ is bounded and $f$ has a maximum or a minimum or both.
    \item $f$ is bounded and $f$ has neither a maximum nor a minimum.
    \item $f$ is unbounded and $f$ has either no maximum or no minimum.
\end{enumerate}

\pyqheader{2026}{7}{MCQ}{1}
Let $f:\mathbb{R}\rightarrow\mathbb{R}$ be a twice differentiable function such that $f(n)=1$ for all $n\in\mathbb{Z}$. Which ONE of the following statements is necessarily TRUE?
\begin{enumerate}[label=(\Alph*)]
    \item $f^{\prime}(n)=0$ for all $n\in\mathbb{Z}$
    \item $f^{\prime}(x)=0$ for infinitely many $x\in\mathbb{R}\setminus\mathbb{Z}$
    \item $f^{\prime\prime}(n)=0$ for all $n\in\mathbb{Z}$
    \item $f^{\prime\prime}(x)=0$ for infinitely many $x\in\mathbb{R}\setminus\mathbb{Z}$
\end{enumerate}

\pyqheader{2026}{13}{MCQ}{2}
Let $f:\mathbb{R}\rightarrow\mathbb{R}$ be a function that has derivatives of all orders. Which ONE of the following statements is FALSE?
\begin{enumerate}[label=(\Alph*)]
    \item If $f$ is a polynomial in $x$ of degree at most 5, then $f^{(6)}(x)=0$ for all $x\in\mathbb{R}$.
    \item If $f^{(6)}(x)=0$ for all $x\in\mathbb{R}$, then $f$ is a polynomial in $x$ of degree at most 5.
    \item If $f^{(k)}(a)=0$ for all $1\le k\le 5$ and $f^{(6)}(a)>0$ for some $a\in\mathbb{R}$, then $f(x)$ has a local minimum at $a$.
    \item If $f^{(k)}(a)=0$ for all $1\le k\le 6$ and $f^{(7)}(a)<0$ for some $a\in\mathbb{R}$, then $f(x)$ has a local maximum at $a$.
\end{enumerate}

\pyqheader{2026}{19}{MCQ}{2}
The value of $\lim_{x\rightarrow\infty}\int_{0}^{x}e^{-t}|\sin t|\,dt$ is:
\begin{enumerate}[label=(\Alph*)]
    \item $\frac{e^{\pi}+1}{2(e^{\pi}-1)}$
    \item $\frac{e^{\pi}-1}{2(e^{\pi}+1)}$
    \item $\frac{2(e^{\pi}+1)}{e^{\pi}-1}$
    \item $\frac{2(e^{\pi}-1)}{e^{\pi}+1}$
\end{enumerate}

\pyqheader{2026}{21}{MCQ}{2}
Let $f:[1,9]\rightarrow\mathbb{R}$ be a non-constant continuous function such that $f(1)=f(9)$. Which ONE of the following statements is necessarily TRUE?
\begin{enumerate}[label=(\Alph*)]
    \item There exists $c\in[4,8]$ such that $f(c)=f(c+1)$.
    \item There exists $c\in[3,7]$ such that $f(c)=f(c+2)$.
    \item There exists $c\in[2,6]$ such that $f(c)=f(c+3)$.
    \item There exists $c\in[1,5]$ such that $f(c)=f(c+4)$.
\end{enumerate}

\pyqheader{2026}{31}{MSQ}{2}
Which of the following functions $f:\mathbb{R}\rightarrow\mathbb{R}$ has/have a local minimum at $x=0$?
\begin{enumerate}[label=(\Alph*)]
    \item $f(x)=\sin|x|$
    \item $f(x)=\sin x+\frac{x^{3}}{6}$
    \item $f(x)=x^{4}+x^{2}+3$
    \item $f(x)=\min\{x-[x],1-x+[x]\}$
\end{enumerate}

\pyqheader{2026}{33}{MSQ}{2}
Which of the following functions is/are differentiable at $x=1$?
\begin{enumerate}[label=(\Alph*)]
    \item $f(x)=|x-1|^{3}$
    \item $f(x)=|x^{2}-1|$
    \item $f(x)=\begin{cases}x^{2}e^{-x^{2}},& |x|\le 1\\ e^{-1},& |x|>1\end{cases}$
    \item $f(x)=[x]$
\end{enumerate}

\pyqheader{2026}{34}{MSQ}{2}
Let $f:\mathbb{R}\rightarrow\mathbb{R}$ be given by $f(x)=\begin{cases}\cos x,& x\in\mathbb{Q}\\ 0,& x\notin\mathbb{Q}\end{cases}$. Which of the following is/are TRUE?
\begin{enumerate}[label=(\Alph*)]
    \item $f(x)$ is continuous at 0.
    \item $f(x)$ is continuous at $\frac{\pi}{2}$.
    \item $f(x)$ is Riemann integrable on $[0,1]$ and $\int_{0}^{1}f(x)dx=\sin 1$.
    \item $f(x)$ is Riemann integrable on $[0,1]$ and $\int_{0}^{1}f(x)dx=0$.
\end{enumerate}

\pyqheader{2026}{44}{NAT}{1}
$\lim_{x\rightarrow 0}\left(\frac{1}{x}-\frac{1}{\sin x}+e^{\frac{1-\cos x}{x}}\right)=$ \underline{\hspace{1.5cm}} (rounded off to one decimal place).

\pyqheader{2026}{48}{NAT}{1}
$\int_{0}^{3}(|x-1|-x[x])\,dx=$ \underline{\hspace{1.5cm}} (rounded off to one decimal place).

\pyqheader{2026}{54}{NAT}{2}
Let $f(x)=\int_{0}^{e^{\sin x}}\int_{0}^{\log_{e}y}e^{-(\frac{1}{\sqrt{2}})}(1-t^{2})dt\,dy$. Then $f^{\prime}(\pi/4)=$ \underline{\hspace{1.5cm}} (rounded off to one decimal place).

\pyqheader{2025}{3}{MCQ}{1}
Let $f,g:\mathbb{R}\rightarrow\mathbb{R}$ be defined by $f(x)=\begin{cases}x|x|\sin\frac{1}{x},& x\ne 0\\ 0,& x=0\end{cases}$ and $g(x)=\begin{cases}x^{2}\sin\frac{1}{x}+x\cos\frac{1}{x},& x\ne 0\\ 0,& x=0\end{cases}$. Then:
\begin{enumerate}[label=(\Alph*)]
    \item $f$ is differentiable at $x=0$, and $g$ is NOT differentiable at $x=0$
    \item $f$ is NOT differentiable at $x=0$, and $g$ is differentiable at $x=0$
    \item $f$ is differentiable at $x=0$ and $g$ is differentiable at $x=0$
    \item $f$ is NOT differentiable at $x=0$, and $g$ is NOT differentiable at $x=0$
\end{enumerate}

\pyqheader{2025}{4}{MCQ}{1}
Let $f,g:\mathbb{R}\rightarrow\mathbb{R}$ be defined by $f(x)=\begin{cases}|x|^{1/3}|\sin\frac{1}{x}|\cos x,& x\ne 0\\ 0,& x=0\end{cases}$ and $g(x)=\begin{cases}e^{x}\cos\frac{1}{x},& x\ne 0\\ 1,& x=0\end{cases}$. Then:
\begin{enumerate}[label=(\Alph*)]
    \item $f$ is continuous at $x=0$, and $g$ is NOT continuous at $x=0$
    \item $f$ is NOT continuous at $x=0$, and $g$ is continuous at $x=0$
    \item $f$ is continuous at $x=0$, and $g$ is continuous at $x=0$
    \item $f$ is NOT continuous at $x=0$, and $g$ is NOT continuous at $x=0$
\end{enumerate}

\pyqheader{2025}{17}{MCQ}{2}
Let $f(x)=10x^{2}+e^{x}-\sin(2x)-\cos x, x\in\mathbb{R}$. The number of points at which the function $f$ has a local minimum is:
\begin{enumerate}[label=(\Alph*)]
    \item 0
    \item 1
    \item 2
    \item greater than or equal to 3
\end{enumerate}

\pyqheader{2025}{30}{MCQ}{2}
The number of elements in the set $\{x\in\mathbb{R}: 8x^{2}+x^{4}+x^{8}=\cos x\}$ is equal to:
\begin{enumerate}[label=(\Alph*)]
    \item 0
    \item 1
    \item 2
    \item greater than or equal to 3
\end{enumerate}

\pyqheader{2025}{37}{MSQ}{2}
Let $f:\mathbb{R}\rightarrow\mathbb{R}$ be a twice differentiable function such that $f(0)=0, f^{\prime}(0)=2$ and $f(1)=-3$. Then, which of the following is/are TRUE?
\begin{enumerate}[label=(\Alph*)]
    \item $|f^{\prime}(x)|\le 2$ for all $x\in[0,1]$
    \item $|f^{\prime}(x_{1})|>2$ for some $x_{1}\in[0,1]$
    \item $|f^{\prime\prime}(x)|<10$ for all $x\in[0,1]$
    \item $|f^{\prime\prime}(x_{2})|\ge 10$ for some $x_{2}\in[0,1]$
\end{enumerate}

\pyqheader{2025}{38}{MSQ}{2}
Let $f:\mathbb{R}\rightarrow\mathbb{R}$ be a twice differentiable function such that $f(0)=4, f(1)=-2, f(2)=8$ and $f(3)=2$. Then, which of the following is/are TRUE?
\begin{enumerate}[label=(\Alph*)]
    \item $|f^{\prime}(x)|<5$ for all $x\in[0,1]$
    \item $|f^{\prime}(x_{1})|\ge 5$ for some $x_{1}\in[0,1]$
    \item $f^{\prime}(x_{2})=0$ for some $x_{2}\in[0,3]$
    \item $f^{\prime\prime}(x_{3})=0$ for some $x_{3}\in[0,3]$
\end{enumerate}

\pyqheader{2025}{42}{NAT}{1}
$\lim_{n\rightarrow\infty}8n\left(e^{(\frac{1}{2n})}-1\right)\left(\sin\frac{1}{2n}+\left|\cos\frac{1}{2n}\right|\right)=$ \underline{\hspace{1.5cm}} (rounded off to two decimal places).

\pyqheader{2025}{43}{NAT}{1}
Let $\alpha$ be the real number such that $\lim_{x\rightarrow 0}\frac{(1-\cos x)(2^{2+x}-4)}{x^{3}}=\alpha\ln 2$. Then, $\alpha=$ \underline{\hspace{1.5cm}} (rounded off to two decimal places).

\pyqheader{2025}{51}{NAT}{2}
Let $f:\mathbb{R}\rightarrow\mathbb{R}$ be continuous satisfying $\int_{0}^{\pi/4}\left(\sin(x)f(x)+\cos(x)\int_{0}^{x}f(t)dt\right)dx=\sqrt{2}$. Then, $\int_{0}^{\pi/4}f(x)\,dx=$ \underline{\hspace{1.5cm}} (rounded off to two decimal places).

\pyqheader{2025}{53}{NAT}{2}
Let $f(x)=2x-\sin x$ for $x\in\mathbb{R}$. If $\lim_{x\rightarrow 0}\left(\frac{1}{x}\sum_{i=1}^{k}i^{2}f(\frac{x}{i})\right)=45$, then $k=$ \underline{\hspace{1.5cm}}.

\pyqheader{2024}{19}{MCQ}{2}
Let $f:[-2\pi,2\pi]\rightarrow\mathbb{R}$ be defined by $f(x)=\sin(\min\{x, x-[x]\})$. Consider the set $S=\{x\in[-2\pi,2\pi]: f \text{ is discontinuous at } x\}$. Then:
\begin{enumerate}[label=(\Alph*)]
    \item $S$ has 13 elements
    \item $S$ has 7 elements
    \item $S$ is infinite
    \item $S$ has 6 elements
\end{enumerate}

\pyqheader{2024}{24}{MCQ}{2}
Consider the statements:
\begin{itemize}
    \item[P:] $\exists f,g$ such that $f$ is continuous at $x=1$, $g$ is discontinuous at $x=1$, but $g\circ f$ is continuous at $x=1$.
    \item[Q:] $\exists f,g$ such that both $f,g$ are discontinuous at $x=1$, but $g\circ f$ is continuous at $x=1$.
\end{itemize}
\begin{enumerate}[label=(\Alph*)]
    \item Both P and Q are true
    \item Both false
    \item P true, Q false
    \item P false, Q true
\end{enumerate}

\pyqheader{2024}{25}{MCQ}{2}
Let $f(x)=\frac{(x^{2}+1)^{2}}{x^{4}+x^{2}+1}$ for $x\in\mathbb{R}$. Then $f$ has:
\begin{enumerate}[label=(\Alph*)]
    \item exactly 2 local maxima and 3 local minima
    \item exactly 3 local maxima and 2 local minima
    \item exactly 1 local maximum and 2 local minima
    \item exactly 2 local maxima and 1 local minimum
\end{enumerate}

\pyqheader{2024}{29}{MCQ}{2}
Let $L_{1}: y=3x+2$ and $L_{2}: y=4x+3$. Suppose $f:\mathbb{R}\rightarrow\mathbb{R}$ is $C^{4}$ such that $L_{1}$ intersects $y=f(x)$ at 3 distinct points and $L_{2}$ intersects $y=f(x)$ at 4 distinct points. Then:
\begin{enumerate}[label=(\Alph*)]
    \item $\frac{df}{dx}$ does not attain 3 on $\mathbb{R}$
    \item $\frac{d^{2}f}{dx^{2}}$ vanishes at most once
    \item $\frac{d^{3}f}{dx^{3}}$ vanishes at least once
    \item $\frac{df}{dx}$ does not attain $\frac{7}{2}$
\end{enumerate}

\pyqheader{2024}{34}{MSQ}{2}
Let $f:(1,\infty)\rightarrow(0,\infty)$ be a continuous function such that $f(n)=n!$ for all $n\in\mathbb{N}$. Then:
\begin{enumerate}[label=(\Alph*)]
    \item $\lim_{x\rightarrow\infty}f(x)$ does not exist
    \item $\lim_{x\rightarrow\infty}f(x)=\infty$
    \item The set of solutions to $f(x)=2024$ is finite
    \item The set of solutions to $f(x)=2024$ is infinite
\end{enumerate}

\pyqheader{2024}{38}{MSQ}{2}
Define $f(x)=\sum_{m=0}^{\infty}\frac{(-1)^{m}x^{2m}}{2^{2m}(m!)^{2}}$ and $g(x)=\frac{x}{2}\sum_{m=0}^{\infty}\frac{(-1)^{m}x^{2m}}{2^{2m}(m+1)!m!}$. Let $0<x_{1}<x_{2}, 0<x_{3}<x_{4}$ with $f(x_{1})=f(x_{2})=0$, $f(x)\ne 0$ on $(x_1,x_2)$, and $g(x_3)=g(x_4)=0$, $g(x)\ne 0$ on $(x_3,x_4)$. Then:
\begin{enumerate}[label=(\Alph*)]
    \item $f$ does not vanish in $(x_3,x_4)$
    \item $f$ vanishes exactly once in $(x_3,x_4)$
    \item $g$ does not vanish in $(x_1,x_2)$
    \item $g$ vanishes exactly once in $(x_1,x_2)$
\end{enumerate}

\pyqheader{2024}{46}{NAT}{1}
Let $S=\{f:\mathbb{R}\rightarrow\mathbb{R}: f \text{ is a polynomial and } f(f(x))=(f(x))^{2024} \text{ for } x\in\mathbb{R}\}$. The number of elements in $S$ is \underline{\hspace{1.5cm}}.

\pyqheader{2024}{52}{NAT}{2}
$\lim_{t\rightarrow\infty}\left(\left(\log(t^{2}+\frac{1}{t^{2}})\right)^{-1}\int_{1}^{\pi t}\frac{\sin^{2}5x}{x}dx\right)=$ \underline{\hspace{1.5cm}} (rounded off to two decimal places).

\pyqheader{2024}{56}{NAT}{2}
Let $f(x)=x^{3}-4x^{2}+4x-6$ and $S(c)=\{x\in\mathbb{R}: f(x)=c\}$. Then $|S(-7)|+|S(-5)|+|S(3)|=$ \underline{\hspace{1.5cm}}.

\pyqheader{2024}{60}{NAT}{2}
Let $f(x)=\sin^{-1}x$. If $a_{6}$ denotes the coefficient of $x^{6}$ in the Taylor series of $(f(x))^{2}$ about $x=0$, then $9a_{6}=$ \underline{\hspace{1.5cm}} (rounded off to two decimal places).

\pyqheader{2023}{5}{MCQ}{1}
Let $p(x)=x^{57}+3x^{10}-21x^{3}+x^{2}+21$ and $q(x)=p(x)+\sum_{j=1}^{57}p^{(j)}(x)$ for all $x\in\mathbb{R}$. Then $q$ admits:
\begin{enumerate}[label=(\Alph*)]
    \item NEITHER a global maximum NOR a global minimum on $\mathbb{R}$
    \item a global maximum but NOT a global minimum on $\mathbb{R}$
    \item a global minimum but NOT a global maximum on $\mathbb{R}$
    \item a global minimum and a global maximum on $\mathbb{R}$
\end{enumerate}

\pyqheader{2023}{6}{MCQ}{1}
The limit $\lim_{x\rightarrow 0}\frac{\int_{0}^{x}\sin(t^{2})dt}{\int_{0}^{x}(\ln(t+1))^{2}dt}$ is:
\begin{enumerate}[label=(\Alph*)]
    \item 0
    \item 1
    \item $\frac{\pi}{e}$
    \item non-existent
\end{enumerate}

\pyqheader{2023}{25}{MCQ}{2}
Let $f(x)=\cos(x)$ and $g(x)=1-\frac{x^{2}}{2}$ for $x\in(-\frac{\pi}{2},\frac{\pi}{2})$. Then:
\begin{enumerate}[label=(\Alph*)]
    \item $f(x)\ge g(x)$ for all $x\in(-\frac{\pi}{2},\frac{\pi}{2})$
    \item $f(x)\le g(x)$ for all $x\in(-\frac{\pi}{2},\frac{\pi}{2})$
    \item $f(x)-g(x)$ changes sign exactly once on $(-\frac{\pi}{2},\frac{\pi}{2})$
    \item $f(x)-g(x)$ changes sign more than once on $(-\frac{\pi}{2},\frac{\pi}{2})$
\end{enumerate}

\pyqheader{2023}{28}{MCQ}{2}
Let $y:\mathbb{R}\rightarrow\mathbb{R}$ be $C^2$ such that $y(0)=y(1)=0$ and $y^{\prime\prime}(x)+x^{2}<0$ for all $x\in[0,1]$. Then:
\begin{enumerate}[label=(\Alph*)]
    \item $y(x)>0$ for all $x\in(0,1)$
    \item $y(x)<0$ for all $x\in(0,1)$
    \item $y(x)=0$ has exactly one solution in $(0,1)$
    \item $y(x)=0$ has more than one solution in $(0,1)$
\end{enumerate}

\pyqheader{2023}{30}{MCQ}{2}
Let $f:\mathbb{R}\rightarrow\mathbb{R}$ be an infinitely differentiable function such that $f^{\prime\prime}$ has exactly two distinct zeroes. Then:
\begin{enumerate}[label=(\Alph*)]
    \item $f^{\prime}$ has at most 3 distinct zeroes
    \item $f^{\prime}$ has at least 1 zero
    \item $f$ has at most 3 distinct zeroes
    \item $f$ has at least 2 distinct zeroes
\end{enumerate}

\pyqheader{2023}{34}{MSQ}{2}
Let $f:(-1,1)\rightarrow\mathbb{R}$ be differentiable satisfying $f(0)=0$ and $|f^{\prime}(x)|\le M|x|$ for all $x\in(-1,1)$. Then:
\begin{enumerate}[label=(\Alph*)]
    \item $f^{\prime}$ is continuous at $x=0$
    \item $f^{\prime}$ is differentiable at $x=0$
    \item $ff^{\prime}$ is differentiable at $x=0$
    \item $(f^{\prime})^{2}$ is differentiable at $x=0$
\end{enumerate}

\pyqheader{2023}{35}{MSQ}{2}
Which of the following functions is/are Riemann integrable on $[0,1]$?
\begin{enumerate}[label=(\Alph*)]
    \item $f(x)=\int_{0}^{x}|\frac{1}{2}-t|dt$
    \item $f(x)=\begin{cases}x\sin(1/x),& \text{if } x\ne 0\\ 0,& \text{if } x=0\end{cases}$
    \item $f(x)=\begin{cases}1,& \text{if } x\in\mathbb{Q}\cap[0,1]\\ -1,& \text{otherwise}\end{cases}$
    \item $f(x)=\begin{cases}x,& \text{if } x\in[0,1)\\ 0,& \text{if } x=1\end{cases}$
\end{enumerate}

\pyqheader{2023}{44}{NAT}{1}
$\lim_{n\rightarrow\infty}\left(n\int_{0}^{1}\frac{x^{n}}{x+1}dx\right)=$ \underline{\hspace{1.5cm}} (rounded off to two decimal places).

\pyqheader{2023}{47}{NAT}{1}
Let $f(x)=\sqrt[3]{x}$ for $x>0$, and $\theta(h)$ satisfy $f(3+h)-f(3)=hf^{\prime}(3+\theta(h)h)$ for $h\in(-1,1)$. Then $\lim_{h\rightarrow 0}\theta(h)=$ \underline{\hspace{1.5cm}} (rounded off to two decimal places).

\pyqheader{2023}{52}{NAT}{2}
The global minimum value of $f(x)=|x-1|+|x-2|^{2}$ on $\mathbb{R}$ is equal to \underline{\hspace{1.5cm}} (rounded off to two decimal places).

\pyqheader{2023}{60}{NAT}{2}
Let $f:\mathbb{R}\rightarrow\mathbb{R}$ be bijective with $f(x)=\sum_{n=1}^{\infty}a_{n}x^{n}$ and $f^{-1}(x)=\sum_{n=1}^{\infty}b_{n}x^{n}$. If $a_{1}=2$ and $a_{2}=4$, then $b_{1}=$ \underline{\hspace{1.5cm}}.

\pyqheader{2022}{22}{MCQ}{2}
Let $H(x)=\frac{1}{2}(e^{x}+e^{-x})$ and $f(x)=\int_{0}^{\pi}H(x\sin\theta)d\theta$. Then:
\begin{enumerate}[label=(\Alph*)]
    \item $xf^{\prime\prime}(x)+f^{\prime}(x)+xf(x)=0$
    \item $xf^{\prime\prime}(x)-f^{\prime}(x)+xf(x)=0$
    \item $xf^{\prime\prime}(x)+f^{\prime}(x)-xf(x)=0$
    \item $xf^{\prime\prime}(x)-f^{\prime}(x)-xf(x)=0$
\end{enumerate}

\pyqheader{2022}{27}{MCQ}{2}
Given $\lim_{N\rightarrow\infty}\int_{0}^{N}e^{-t^{2}}dt=\frac{\sqrt{\pi}}{2}$, the value of $\lim_{N\rightarrow\infty}\int_{0}^{N}\frac{1}{t^{2}}(e^{-at^{2}}-e^{-bt^{2}})dt$ ($0<a<b$) is:
\begin{enumerate}[label=(\Alph*)]
    \item $\sqrt{\pi}(\sqrt{a}-\sqrt{b})$
    \item $\sqrt{\pi}(\sqrt{a}+\sqrt{b})$
    \item $-\sqrt{\pi}(\sqrt{a}+\sqrt{b})$
    \item $\sqrt{\pi}(\sqrt{b}-\sqrt{a})$
\end{enumerate}

\pyqheader{2022}{29}{MCQ}{2}
For $n\in\mathbb{N}$ and $x\ge 1$, let $f_{n}(x)=\int_{0}^{\pi}(x^{2}+(\cos\theta)\sqrt{x^{2}-1})^{n}d\theta$. Then:
\begin{enumerate}[label=(\Alph*)]
    \item $f_{n}(x)$ is not a polynomial in $x$ if $n$ is odd and $n\ge 3$.
    \item $f_{n}(x)$ is not a polynomial in $x$ if $n$ is even and $n\ge 4$.
    \item $f_{n}(x)$ is a polynomial in $x$ for all $n\in\mathbb{N}$.
    \item $f_{n}(x)$ is not a polynomial in $x$ for any $n\ge 3$.
\end{enumerate}

\pyqheader{2022}{32}{MSQ}{2}
Let $S$ be the set of continuous $f:[-1,1]\rightarrow\mathbb{R}$ that are analytic on $(-1,1)$ with Taylor series converging to $f(x)$, and $f(1/n)=0$ for all $n\in\mathbb{N}$. Then:
\begin{enumerate}[label=(\Alph*)]
    \item $f(0)=0$ for every $f\in S$.
    \item $f^{\prime}(1/2)=0$ for every $f\in S$.
    \item There exists $f\in S$ such that $f^{\prime}(1/2)\ne 0$.
    \item There exists $f\in S$ such that $f(x)\ne 0$ for some $x\in[-1,1]$.
\end{enumerate}

\pyqheader{2022}{33}{MSQ}{2}
Define $f:[0,1]\rightarrow[0,1]$ by $f(x)=\begin{cases}1/2,& x=0\\ 1/n,& x=m/n \text{ in lowest terms}\\ 0,& x\text{ irrational}\end{cases}$ and $g(x)=\begin{cases}0,& x=0\\ 1,& x\in(0,1]\end{cases}$. Then:
\begin{enumerate}[label=(\Alph*)]
    \item $f$ is Riemann integrable on $[0,1]$.
    \item $g$ is Riemann integrable on $[0,1]$.
    \item $f\circ g$ is Riemann integrable on $[0,1]$.
    \item $g\circ f$ is Riemann integrable on $[0,1]$.
\end{enumerate}

\pyqheader{2022}{34}{MSQ}{2}
Let $S=\{f:\mathbb{R}\rightarrow\mathbb{R} : |f(x)-f(y)|^{2}\le |x-y|^{3} \text{ for all } x,y\in\mathbb{R}\}$. Then:
\begin{enumerate}[label=(\Alph*)]
    \item Every function in $S$ is differentiable.
    \item There exists $f\in S$ such that $f$ is differentiable, but not twice differentiable.
    \item There exists $f\in S$ such that $f$ is twice differentiable, but not thrice differentiable.
    \item Every function in $S$ is infinitely differentiable.
\end{enumerate}

\pyqheader{2022}{41}{NAT}{1}
$\lim_{n\rightarrow\infty}\left(\frac{1^{4}+2^{4}+\cdots+n^{4}}{n^{5}}+\frac{1}{\sqrt{n}}\left(\frac{1}{\sqrt{n+1}}+\cdots+\frac{1}{\sqrt{4n}}\right)\right)=$ \underline{\hspace{1.5cm}} (rounded off to two decimal places).

\pyqheader{2022}{54}{NAT}{2}
Let $f(x)=\begin{cases}(x-\pi)e^{\sin x},& 0\le x\le\pi/2\\ xe^{\sin x}+\frac{4}{\pi},& \pi/2<x\le\pi\end{cases}$. Then $\int_{0}^{\pi}f(x)\,dx=$ \underline{\hspace{1.5cm}} (rounded off to two decimal places).

\pyqheader{2022}{59}{NAT}{2}
Let $f,g:(-1,1)\rightarrow\mathbb{R}$ be $C^3$ with $f(x)\ne g(x)$ for $x\ne 0$, and $f(0)=g(0)=\ln 2, f^{\prime}(0)=g^{\prime}(0)=\pi, f^{\prime\prime}(0)=g^{\prime\prime}(0)=\pi^2, f^{\prime\prime\prime}(0)=\pi^9, g^{\prime\prime\prime}(0)=\pi^3$. Then $\lim_{x\rightarrow 0}\frac{e^{f(x)}-e^{g(x)}}{f(x)-g(x)}=$ \underline{\hspace{1.5cm}} (rounded off to two decimal places).

\pyqheader{2022}{60}{NAT}{2}
If $\int_{0}^{x^{3}+x^{2}}f(t)dt=x^{2}$ and $\int_{0}^{g(x)}t^{2}dt=9(x+1)^{3}$, then $f(2)+g(2)+16f(12)=$ \underline{\hspace{1.5cm}} (rounded off to two decimal places).

% ==============================================================================
% SECTION 2: MULTIVARIABLE CALCULUS & DIFFERENTIAL EQUATIONS
% ==============================================================================
\section{Multivariable Calculus and Differential Equations}

\subsection{Functions of Two or Three Real Variables}

\pyqheader{2026}{15}{MCQ}{2}
Let $f:\mathbb{R}^{2}\rightarrow\mathbb{R}$ be given by $f(x,y)=(x^{2}-1)^{2}+(y^{2}-1)^{2}$. Which ONE of the following is TRUE?
\begin{enumerate}[label=(\Alph*)]
    \item $f$ has local maxima at exactly two points.
    \item $f$ has local minima at exactly two points.
    \item $f$ has local minima at exactly three points.
    \item $f$ has exactly four saddle points.
\end{enumerate}

\pyqheader{2026}{17}{MCQ}{2}
Let $f(x,y)=\begin{cases}\frac{5x^{2}y-4y^{3}}{x^{2}+y^{2}},& (x,y)\ne(0,0)\\ 0,& (x,y)=(0,0)\end{cases}$. Which ONE of the following statements is FALSE?
\begin{enumerate}[label=(\Alph*)]
    \item $f$ is continuous at $(0,0)$
    \item $\frac{\partial f}{\partial x}$ at $(0,1)=0$
    \item $\frac{\partial f}{\partial y}$ at $(1,0)=5$
    \item $\frac{\partial f}{\partial y}$ is continuous at $(0,0)$
\end{enumerate}

\pyqheader{2026}{36}{MSQ}{2}
Let $f(x,y)=\begin{cases}\frac{x^{2}y}{1+x^{2}}\sin(1/x),& x\ne 0\\ 0,& x=0\end{cases}$. Which of the following statements is/are TRUE?
\begin{enumerate}[label=(\Alph*)]
    \item $\lim_{y\rightarrow 0}\lim_{x\rightarrow 0}f(x,y)$ exists.
    \item $\frac{\partial f}{\partial x}$ is continuous at the point $(0,1)$.
    \item $\frac{\partial f}{\partial y}$ is continuous at the point $(0,1)$.
    \item $\lim_{(x,y)\rightarrow(0,0)}\frac{f(x,y)}{\sqrt{x^{2}+y^{2}}}$ does not exist.
\end{enumerate}

\pyqheader{2026}{46}{NAT}{1}
Let $z=\cos(4x+5y)$, where $x=\frac{\pi}{2}+2\theta, y=-(\frac{\pi}{4}+\theta)$. Then $\frac{dz}{d\theta}\Big|_{\theta=\frac{\pi}{4}}=$ \underline{\hspace{1.5cm}} (rounded off to one decimal place).

\pyqheader{2026}{52}{NAT}{2}
Let $f(x,y)=\begin{cases}4x^{2}\tan^{-1}(\frac{y}{2x})+y^{3}\tan^{-1}(\frac{x}{4y^{2}}),& xy\ne 0\\ 0,& \text{otherwise}\end{cases}$. Then $\frac{\partial^{2}f}{\partial y\partial x}(0,0)=$ \underline{\hspace{1.5cm}} (rounded off to two decimal places).

\pyqheader{2025}{27}{MCQ}{2}
Let $f:\mathbb{R}^{2}\rightarrow\mathbb{R}$ be defined by $f(x,y)=e^{y}(x^{2}+y^{2})$. Then:
\begin{enumerate}[label=(\Alph*)]
    \item $f$ has 2 local minima
    \item $f$ has 2 local maxima
    \item $f$ has 1 local minimum
    \item $f$ has 1 local maximum
\end{enumerate}

\pyqheader{2025}{31}{MSQ}{2}
Let $f(x,y)=\begin{cases}\frac{xy^{2}+y^{5}}{x^{2}+y^{4}},& (x,y)\ne(0,0)\\ 0,& (x,y)=(0,0)\end{cases}$. Then:
\begin{enumerate}[label=(\Alph*)]
    \item The iterated limits $\lim_{x\rightarrow 0}\lim_{y\rightarrow 0}f(x,y)$ and $\lim_{y\rightarrow 0}\lim_{x\rightarrow 0}f(x,y)$ exist
    \item Exactly one of $\frac{\partial f}{\partial x}$ and $\frac{\partial f}{\partial y}$ exists at $(0,0)$
    \item Both partial derivatives $\frac{\partial f}{\partial x}$ and $\frac{\partial f}{\partial y}$ exist at $(0,0)$
    \item $f$ is NOT differentiable at $(0,0)$
\end{enumerate}

\pyqheader{2025}{34}{MSQ}{2}
Let $f(x,y)=\begin{cases}\frac{(x^{2}+\sin x)y^{2}}{x^{2}+y^{2}},& (x,y)\ne(0,0)\\ 0,& (x,y)=(0,0)\end{cases}$. Then:
\begin{enumerate}[label=(\Alph*)]
    \item $\lim_{(x,y)\rightarrow(0,0)}f(x,y)=1$
    \item $\lim_{(x,y)\rightarrow(0,0)}f(x,y)=0$
    \item $f$ is differentiable at $(0,0)$
    \item $f$ is NOT differentiable at $(0,0)$
\end{enumerate}

\pyqheader{2025}{50}{NAT}{1}
Let $f(x,y)=\begin{cases}\frac{(x^{2}-y^{2})xy}{x^{2}+y^{2}},& (x,y)\ne(0,0)\\ 0,& (x,y)=(0,0)\end{cases}$. Then $\frac{\partial f}{\partial y}(1,0)-\frac{\partial f}{\partial x}(0,2)=$ \underline{\hspace{1.5cm}} (rounded off to two decimal places).

\pyqheader{2024}{2}{MCQ}{1}
For a $C^2$ function $g:\mathbb{R}\rightarrow\mathbb{R}$, define $u_{g}(x,y)=\frac{1}{y}\int_{-y}^{y}g(x+t)dt$ for $y>0$. Then:
\begin{enumerate}[label=(\Alph*)]
    \item $\frac{\partial^{2}u_{g}}{\partial x^{2}}=\frac{2}{y}\frac{\partial u_{g}}{\partial y}+\frac{\partial^{2}u_{g}}{\partial y^{2}}$
    \item $\frac{\partial^{2}u_{g}}{\partial x^{2}}=\frac{1}{y}\frac{\partial u_{g}}{\partial y}+\frac{\partial^{2}u_{g}}{\partial y^{2}}$
    \item $\frac{\partial^{2}u_{g}}{\partial x^{2}}=\frac{2}{y}\frac{\partial u_{g}}{\partial y}-\frac{\partial^{2}u_{g}}{\partial y^{2}}$
    \item $\frac{\partial^{2}u_{g}}{\partial x^{2}}=\frac{1}{y}\frac{\partial u_{g}}{\partial y}-\frac{\partial^{2}u_{g}}{\partial y^{2}}$
\end{enumerate}

\pyqheader{2024}{30}{MCQ}{2}
Define $f(x,y)=12xy e^{-(2x+3y-2)}$. If $(a,b)$ is the local maximum point, then $f(a,b)$ equals:
\begin{enumerate}[label=(\Alph*)]
    \item 2
    \item 6
    \item 12
    \item 0
\end{enumerate}

\pyqheader{2024}{35}{MSQ}{2}
Let $S=\{(x,y)\in\mathbb{R}^{2}: x>0, y>0\}$ and $f(x,y)=2x^{2}+3y^{2}-\log x-\frac{1}{6}\log y$. Then:
\begin{enumerate}[label=(\Alph*)]
    \item Unique local maximum point in $S$
    \item Unique local minimum point in $S$
    \item For each $(x_{0},y_{0})\in S$, the set $\{(x,y)\in S: f(x,y)=f(x_{0},y_{0})\}$ is bounded
    \item For each $(x_{0},y_{0})\in S$, the set $\{(x,y)\in S: f(x,y)=f(x_{0},y_{0})\}$ is unbounded
\end{enumerate}

\pyqheader{2023}{8}{MCQ}{1}
Let $f:\mathbb{R}^{2}\rightarrow\mathbb{R}^{2}$ be defined by $f(x,y)=(e^{x}\cos y, e^{x}\sin y)$. Then the number of points in $\mathbb{R}^{2}$ that do NOT lie in the range of $f$ is:
\begin{enumerate}[label=(\Alph*)]
    \item 0
    \item 1
    \item 2
    \item infinite
\end{enumerate}

\pyqheader{2023}{16}{MCQ}{2}
Let $f(x,y)=e^{x^{2}+y^{2}}$ for $(x,y)\in\mathbb{R}^{2}$, and $a_{n}$ be the determinant of the matrix $\begin{pmatrix}f_{xx} & f_{xy}\\ f_{yx} & f_{yy}\end{pmatrix}$ evaluated at $(\cos n, \sin n)$. Then the limit $\lim_{n\rightarrow\infty}a_{n}$ is:
\begin{enumerate}[label=(\Alph*)]
    \item non-existent
    \item 0
    \item $6e^{2}$
    \item $12e^{2}$
\end{enumerate}

\pyqheader{2023}{17}{MCQ}{2}
Let $f(x,y)=\ln(1+x^{2}+y^{2})$. Let $P=f_{xx}, Q=f_{xy}, R=f_{yx}, S=f_{yy}$ at $(0,0)$. Then:
\begin{enumerate}[label=(\Alph*)]
    \item $PS-QR>0$ and $P<0$
    \item $PS-QR>0$ and $P>0$
    \item $PS-QR<0$ and $P>0$
    \item $PS-QR<0$ and $P<0$
\end{enumerate}

\pyqheader{2023}{36}{MSQ}{2}
Which of the following subsets of $\mathbb{R}^{2}$ is/are bounded?
\begin{enumerate}[label=(\Alph*)]
    \item $\{(x,y)\in\mathbb{R}^{2}: e^{x^{2}}+y^{2}\le 4\}$
    \item $\{(x,y)\in\mathbb{R}^{2}: |x|+|y|\le 4\}$
    \item $\{(x,y)\in\mathbb{R}^{2}: x^{4}+y^{2}\le 4\}$
    \item $\{(x,y)\in\mathbb{R}^{2}: e^{x^{3}}+y^{2}\le 4\}$
\end{enumerate}

\pyqheader{2023}{37}{MSQ}{2}
Let $f(x,y)=\begin{cases}\frac{x^{4}y^{3}}{x^{6}+y^{6}},& \text{if } (x,y)\ne(0,0)\\ 0,& \text{if } (x,y)=(0,0)\end{cases}$. Then:
\begin{enumerate}[label=(\Alph*)]
    \item $\lim_{t\rightarrow 0}\frac{f(t,t)-f(0,0)}{t}$ exists and equals $\frac{1}{2}$
    \item $\frac{\partial f}{\partial x}\Big|_{(0,0)}$ exists and equals 0
    \item $\frac{\partial f}{\partial y}\Big|_{(0,0)}$ exists and equals 0
    \item $\lim_{t\rightarrow 0}\frac{f(t,2t)-f(0,0)}{t}$ exists and equals $\frac{1}{3}$
\end{enumerate}

\pyqheader{2023}{41}{NAT}{1}
Let $f(x,y)=\begin{cases}(x^{2}-1)^{2}\cos^{2}(\frac{y^{2}}{x^{2}-1}),& x\ne\pm 1\\ 0,& x=\pm 1\end{cases}$. The number of points of discontinuity of $f(x,y)$ is \underline{\hspace{1.5cm}}.

\pyqheader{2023}{51}{NAT}{2}
Let $f:\mathbb{R}^{3}\rightarrow\mathbb{R}$ be defined as $f(x,y,z)=x^{3}+y^{3}+z^{3}$, and let $L:\mathbb{R}^{3}\rightarrow\mathbb{R}$ be the linear map satisfying
\[
\lim_{(x,y,z)\rightarrow(0,0,0)}\frac{f(1+x,1+y,1+z)-f(1,1,1)-L(x,y,z)}{\sqrt{x^{2}+y^{2}+z^{2}}}=0.
\]
Then $L(1,2,4)$ is equal to \underline{\hspace{1.5cm}} (rounded off to two decimal places).

\pyqheader{2022}{9}{MCQ}{1}
Let $u(x,t)=\frac{1}{2c}\int_{x-ct}^{x+ct}e^{s^{2}}ds$ on $\mathbb{R}^{2}$. Then:
\begin{enumerate}[label=(\Alph*)]
    \item $\frac{\partial^{2}u}{\partial t^{2}}=c^{2}\frac{\partial^{2}u}{\partial x^{2}}$
    \item $\frac{\partial u}{\partial t}=c^{2}\frac{\partial^{2}u}{\partial x^{2}}$
    \item $\frac{\partial u}{\partial t}\frac{\partial u}{\partial x}=0$
    \item $\frac{\partial^{2}u}{\partial t\partial x}=0$
\end{enumerate}

\pyqheader{2022}{10}{MCQ}{1}
Let $u(x,y)=x-\frac{x}{x^{2}+y^{2}}$ and $v(x,y)=y+\frac{y}{x^{2}+y^{2}}$. The value of the Jacobian determinant $\left|\begin{smallmatrix}u_x & u_y\\ v_x & v_y\end{smallmatrix}\right|$ at $(\cos\theta,\sin\theta)$ ($\theta\in(\pi/4,\pi/2)$) is:
\begin{enumerate}[label=(\Alph*)]
    \item $4\sin\theta$
    \item $4\cos\theta$
    \item $4\sin^{2}\theta$
    \item $4\cos^{2}\theta$
\end{enumerate}

\pyqheader{2022}{15}{MCQ}{2}
For $t\in\mathbb{R}$, let $[t]$ denote the floor function. Define $h(x,y)=\begin{cases}\frac{-1}{x^{2}-y},& x^{2}\ne y\\ 0,& x^{2}=y\end{cases}$ and $g(x)=\begin{cases}\frac{\sin x}{x},& x\ne 0\\ 0,& x=0\end{cases}$. Which ONE is FALSE?
\begin{enumerate}[label=(\Alph*)]
    \item $\lim_{(x,y)\rightarrow(\sqrt{2},\pi)}\cos(\frac{x^{2}y}{x^{2}+1})=-\frac{1}{2}$
    \item $\lim_{(x,y)\rightarrow(\sqrt{2},2)}e^{h(x,y)}=0$
    \item $\lim_{(x,y)\rightarrow(e,e)}\ln(x^{y-[y]})=e-2$
    \item $\lim_{(x,y)\rightarrow(0,0)}e^{2y}g(x)=1$
\end{enumerate}

\pyqheader{2022}{36}{MSQ}{2}
Let $M$ be a positive real number and let $u,v:\mathbb{R}^{2}\rightarrow\mathbb{R}$ be continuous functions satisfying $\sqrt{(u(x,y))^{2}+(v(x,y))^{2}}\ge M\sqrt{x^{2}+y^{2}}$ for all $(x,y)\in\mathbb{R}^{2}$. Let $F:\mathbb{R}^{2}\rightarrow\mathbb{R}^{2}$ be given by $F(x,y)=(u(x,y),v(x,y))$. Then which of the following is/are true?
\begin{enumerate}[label=(\Alph*)]
    \item $F$ is injective.
    \item If $K$ is open in $\mathbb{R}^{2}$, then $F(K)$ is open in $\mathbb{R}^{2}$.
    \item If $K$ is closed in $\mathbb{R}^{2}$, then $F(K)$ is closed in $\mathbb{R}^{2}$.
    \item If $E$ is closed and bounded in $\mathbb{R}^{2}$, then $F^{-1}(E)$ is closed and bounded in $\mathbb{R}^{2}$.
\end{enumerate}

\pyqheader{2022}{42}{NAT}{1}
Let $u(x_{1},x_{2},x_{3})=x_{1}x_{2}^{4}x_{3}^{2}-x_{1}^{3}x_{3}^{4}-26x_{1}^{2}x_{2}^{2}x_{3}^{3}$. Let $x_{1}\frac{\partial u}{\partial x_{2}}+2x_{2}\frac{\partial u}{\partial x_{3}}$ at $(t,t^2,t^3)$ be $ct^k$. Then $k=$ \underline{\hspace{1.5cm}}.

\pyqheader{2022}{51}{NAT}{2}
Let $D=\{(x,y)\in\mathbb{R}^{2}: x^{2}+y^{2}<4\}$. Define $f(x,y)=[x^{2}+y^{2}]\frac{x^{2}y^{2}}{x^{4}+y^{4}}$ and $g(x,y)=[y^{2}]\frac{xy}{x^{2}+y^{2}}$ with $f(0,0)=g(0,0)=0$. The number of points in $D$ at which both $f$ and $g$ are discontinuous is \underline{\hspace{1.5cm}}.

\pyqheader{2022}{56}{NAT}{2}
Let $f(x,y)=x^{2}+2y^{2}-x$. Let $D=\{(x,y): x^{2}+y^{2}\le 1\}$ and $E=\{(x,y): \frac{x^2}{4}+\frac{y^2}{9}\le 1\}$. The total number of elements in $D_{\max}\cup D_{\min}\cup E_{\max}\cup E_{\min}$ is \underline{\hspace{1.5cm}}.

% ------------------------------------------------------------------------------
\subsection{Integral Calculus}

\pyqheader{2026}{5}{MCQ}{1}
Let $f:\mathbb{R}^{3}\rightarrow\mathbb{R}$ be continuous. Then $\int_{0}^{2}\int_{0}^{x}\int_{0}^{y}f(x,y,z)\,dz\,dy\,dx =$:
\begin{enumerate}[label=(\Alph*)]
    \item $\int_{0}^{2}\int_{0}^{y}\int_{y}^{2}f(x,y,z)\,dx\,dz\,dy$
    \item $\int_{0}^{2}\int_{0}^{x}\int_{z}^{2}f(x,y,z)\,dy\,dz\,dx$
    \item $\int_{0}^{2}\int_{y}^{2}\int_{0}^{x}f(x,y,z)\,dz\,dx\,dy$
    \item $\int_{0}^{2}\int_{z}^{2}\int_{0}^{x}f(x,y,z)\,dy\,dx\,dz$
\end{enumerate}

\pyqheader{2026}{10}{MCQ}{1}
Let $T$ be the triangular region in the plane with vertices $(0,0), (1,0), (1,1)$. Double integral of $f$ over $T$ in polar coordinates equals:
\begin{enumerate}[label=(\Alph*)]
    \item $\int_{0}^{\sqrt{2}}\left(\int_{0}^{\pi/4}f\,d\theta\right)r\,dr + \int_{1}^{\sqrt{2}}\left(\int_{0}^{\cos^{-1}(1/r)}f\,d\theta\right)r\,dr$
    \item $\int_{0}^{\sqrt{2}}\left(\int_{0}^{\pi/4}f\,d\theta\right)r\,dr - \int_{1}^{\sqrt{2}}\left(\int_{0}^{\cos^{-1}(1/r)}f\,d\theta\right)r\,dr$
    \item $\int_{0}^{\sqrt{2}}\left(\int_{0}^{\pi/4}f\,d\theta\right)r\,dr + \int_{1}^{\sqrt{2}}\left(\int_{0}^{\sin^{-1}(1/r)}f\,d\theta\right)r\,dr$
    \item $\int_{0}^{\sqrt{2}}\left(\int_{0}^{\pi/4}f\,d\theta\right)r\,dr - \int_{1}^{\sqrt{2}}\left(\int_{0}^{\sin^{-1}(1/r)}f\,d\theta\right)r\,dr$
\end{enumerate}

\pyqheader{2026}{20}{MCQ}{2}
The semi-circle $y=\sqrt{6x-5-x^{2}}$ is rotated clockwise by $\frac{\pi}{2}$ in the plane about the origin. The area of the planar region traced out is:
\begin{enumerate}[label=(\Alph*)]
    \item $6\pi$
    \item $4\pi$
    \item $2\pi$
    \item $\pi$
\end{enumerate}

\pyqheader{2026}{22}{MCQ}{2}
Let $S$ be the solid intersection of $S_{1}: x^{2}+y^{2}+(z-1)^{2}\le 4$ and $S_{2}: x^{2}+y^{2}+(z+1)^{2}\le 4$. The volume of $S$ is:
\begin{enumerate}[label=(\Alph*)]
    \item $8\int_{0}^{\sqrt{3}}\int_{0}^{\sqrt{3-x^{2}}}(\sqrt{4-x^{2}-y^{2}}-1)\,dy\,dx$
    \item $8\int_{0}^{2}\int_{0}^{\sqrt{4-x^{2}}}(\sqrt{4-x^{2}-y^{2}}-1)\,dy\,dx$
    \item $2\int_{0}^{2}\int_{0}^{\sqrt{4-x^{2}}}(\sqrt{4-x^{2}-y^{2}}-1)\,dy\,dx$
    \item $2\int_{0}^{\sqrt{3}}\int_{0}^{\sqrt{3-x^{2}}}(\sqrt{4-x^{2}-y^{2}}-1)\,dy\,dx$
\end{enumerate}

\pyqheader{2026}{53}{NAT}{2}
The double integral of $f(x,y)=x$ over the triangular region with vertices $(-\frac{1}{2},\frac{1}{2}), (1,2)$ and $(1,-1)$ is \underline{\hspace{1.5cm}} (rounded off to one decimal place).

\pyqheader{2026}{55}{NAT}{2}
The volume of the tetrahedron bounded by $x=1, y=2, z=3$ and $12x+8y+6z=70$ is \underline{\hspace{1.5cm}} (rounded off to one decimal place).

\pyqheader{2025}{10}{MCQ}{1}
The value of $\int_{0}^{1}\left(\int_{\sqrt{y}}^{1}3e^{x^{3}}dx\right)dy$ is equal to:
\begin{enumerate}[label=(\Alph*)]
    \item $e-1$
    \item $\frac{e-1}{2}$
    \item $\sqrt{e}-1$
    \item $\frac{\sqrt{e}-1}{2}$
\end{enumerate}

\pyqheader{2025}{28}{MCQ}{2}
Let $\Omega$ be the bounded region in $\mathbb{R}^{3}$ in the first octant bounded by $z=x^{2}+y^{2}, z=4, x=0, y=0$. The volume of $\Omega$ is:
\begin{enumerate}[label=(\Alph*)]
    \item $\pi$
    \item $2\pi$
    \item $3\pi$
    \item $4\pi$
\end{enumerate}

\pyqheader{2025}{45}{NAT}{1}
Let $S$ be the surface area of the portion of the plane $z=x+y+3$ inside $x^{2}+y^{2}=1$. Then $(\frac{S}{\pi})^{2}=$ \underline{\hspace{1.5cm}} (rounded off to two decimal places).

\pyqheader{2025}{48}{NAT}{1}
Let $T$ denote the triangle bounded by the $x$-axis, $y=x$ and $x=1$. Then $\iint_{T}(5-y)dx\,dy=$ \underline{\hspace{1.5cm}} (rounded off to two decimal places).

\pyqheader{2025}{57}{NAT}{2}
Let $\Omega$ be the solid bounded by $z=0, y=0, x=\frac{1}{2}, 2y=x$ and $2x+y+z=4$. If $V$ is the volume, then $64V=$ \underline{\hspace{1.5cm}}.

\pyqheader{2024}{27}{MCQ}{2}
For $a>b>0$, let $D=\{(x,y,z)\in\mathbb{R}^{3}: x^{2}+y^{2}+z^{2}\le a^{2} \text{ and } x^{2}+y^{2}\ge b^{2}\}$. The surface area of the boundary of $D$ is:
\begin{enumerate}[label=(\Alph*)]
    \item $4\pi(a+b)\sqrt{a^{2}-b^{2}}$
    \item $4\pi(a^{2}-b\sqrt{a^{2}-b^{2}})$
    \item $4\pi(a-b)\sqrt{a^{2}-b^{2}}$
    \item $4\pi(a^{2}+b\sqrt{a^{2}-b^{2}})$
\end{enumerate}

\pyqheader{2024}{41}{NAT}{1}
The area of the region $R=\{(x,y)\in\mathbb{R}^{2}: 0\le x\le 1, 0\le y\le 1, \frac{1}{4}\le xy\le\frac{1}{2}\}$ is \underline{\hspace{1.5cm}} (rounded off to two decimal places).

\pyqheader{2024}{48}{NAT}{1}
Let $S=\{(x,y,z)\in\mathbb{R}^{3}: x^{2}+y^{2}+z^{2}=4, (x-1)^{2}+y^{2}\le 1, z\ge 0\}$. The surface area of $S$ equals \underline{\hspace{1.5cm}} (rounded off to two decimal places).

\pyqheader{2024}{53}{NAT}{2}
Let $T$ be the square region enclosed by $(0,1), (1,0), (0,-1), (-1,0)$. Then $\iint_{T}(\cos(\pi(x-y))-\cos(\pi(x+y)))^{2}dx\,dy=$ \underline{\hspace{1.5cm}} (rounded off to two decimal places).

\pyqheader{2024}{54}{NAT}{2}
Let $S=\{(x,y,z)\in\mathbb{R}^{3}: x^{2}+y^{2}+z^{2}<1\}$. Then $\frac{1}{\pi}\iiint_{S}((x-2y+z)^{2}+(2x-y-z)^{2}+(x-y+2z)^{2})dx\,dy\,dz=$ \underline{\hspace{1.5cm}} (rounded off to two decimal places).

\pyqheader{2024}{57}{NAT}{2}
Let $c>0$ be such that $\int_{0}^{c}e^{s^{2}}ds=3$. Then $\int_{0}^{c}\left(\int_{x}^{c}e^{x^{2}+y^{2}}dy\right)dx=$ \underline{\hspace{1.5cm}} (rounded off to one decimal place).

\pyqheader{2023}{7}{MCQ}{1}
$\int_{0}^{1}\int_{0}^{1-x}\cos(x^{3}+y^{2})dy\,dx - \int_{0}^{1}\int_{0}^{1-y}\cos(x^{3}+y^{2})dx\,dy=$:
\begin{enumerate}[label=(\Alph*)]
    \item 0
    \item $\frac{\cos 1}{2}$
    \item $\frac{\sin 1}{2}$
    \item $\cos(\frac{1}{2})-\sin(\frac{1}{2})$
\end{enumerate}

\pyqheader{2023}{18}{MCQ}{2}
The curved surface area of $S=\{(x,y,z)\in\mathbb{R}^{3}: z^{2}=(x-1)^{2}+(y-2)^{2}\}$ between $z=2$ and $z=3$ is:
\begin{enumerate}[label=(\Alph*)]
    \item $4\pi\sqrt{2}$
    \item $5\pi\sqrt{2}$
    \item $9\pi$
    \item $9\pi\sqrt{2}$
\end{enumerate}

\pyqheader{2023}{26}{MCQ}{2}
Let $f(x,y)=\iint_{(u-x)^{2}+(v-y)^{2}\le 1}e^{-\sqrt{(u-x)^{2}+(v-y)^{2}}}du\,dv$. Then $\lim_{n\rightarrow\infty}f(n,n^{2})$ is:
\begin{enumerate}[label=(\Alph*)]
    \item non-existent
    \item 0
    \item $\pi(1-e^{-1})$
    \item $2\pi(1-2e^{-1})$
\end{enumerate}

\pyqheader{2023}{31}{MSQ}{2}
For $t\in(0,1)$, let $P_{t}=\{(x,y,z): (x^{2}+y^{2})z=1, t^{2}\le x^{2}+y^{2}\le 1\}$ and $a_{t}$ be its surface area. Then:
\begin{enumerate}[label=(\Alph*)]
    \item $a_{t}=\iint_{t^{2}\le x^{2}+y^{2}\le 1}\sqrt{1+\frac{4x^{2}}{(x^{2}+y^{2})^{4}}+\frac{4y^{2}}{(x^{2}+y^{2})^{4}}}\,dx\,dy$
    \item $a_{t}=\iint_{t^{2}\le x^{2}+y^{2}\le 1}\sqrt{1+\frac{4x^{2}}{(x^{2}+y^{2})^{2}}+\frac{4y^{2}}{(x^{2}+y^{2})^{2}}}\,dx\,dy$
    \item $\lim_{t\rightarrow 0^{+}}a_{t}$ does NOT exist
    \item $\lim_{t\rightarrow 0^{+}}a_{t}$ exists
\end{enumerate}

\pyqheader{2023}{48}{NAT}{1}
Let $V$ be the volume of $S=\{(x,y,z): xy\le z\le 4, x^{2}+y^{2}\le 1\}$. Then $\frac{V}{\pi}=$ \underline{\hspace{1.5cm}} (rounded off to two decimal places).

\pyqheader{2023}{55}{NAT}{2}
Let $S$ be the triangular region with vertices $(0,0), (0,\frac{\pi}{2})$ and $(\frac{\pi}{2},0)$. Then $\iint_{S}\sin x\cos y\,dx\,dy=$ \underline{\hspace{1.5cm}} (rounded off to two decimal places).

\pyqheader{2022}{11}{MCQ}{2}
Let $G=\{(s,t): 0<s<1, 0<t<1\}$ and $T(s,t)=(\frac{\pi s(1-t)}{2},\frac{\pi(1-s)}{2})$. The area of $T(G)$ is:
\begin{enumerate}[label=(\Alph*)]
    \item $\frac{\pi}{4}$
    \item $\frac{\pi^{2}}{4}$
    \item $\frac{\pi^{2}}{8}$
    \item 1
\end{enumerate}

\pyqheader{2022}{48}{NAT}{1}
Consider the region $G=\{(x,y,z)\in\mathbb{R}^{3}: 0<z<x^{2}-y^{2}, x^{2}+y^{2}<1\}$. Then the volume of $G$ is equal to \underline{\hspace{1.5cm}} (rounded off to two decimal places).

\pyqheader{2022}{52}{NAT}{2}
If $G=\{(x,y): x^{2}+y^{2}<1, \frac{x}{\sqrt{3}}<y<\sqrt{3}x, x>0, y>0\}$, then $\frac{200}{\pi}\iint_{G}x^{2}dx\,dy=$ \underline{\hspace{1.5cm}} (rounded off to two decimal places).

% ------------------------------------------------------------------------------
\subsection{Differential Equations}

\pyqheader{2026}{4}{MCQ}{1}
The general solution of $x\sin(\frac{y}{x})\frac{dy}{dx}=y\sin(\frac{y}{x})+\frac{x}{2}$ ($x>0$) is:
\begin{enumerate}[label=(\Alph*)]
    \item $\cos(\frac{x}{y})+\log_{e}(x^{2})=K$
    \item $\cos(\frac{y}{x})+\log_{e}(x^{2})=K$
    \item $\cos(\frac{x}{y})+\log_{e}(\sqrt{x})=K$
    \item $\cos(\frac{y}{x})+\log_{e}(\sqrt{x})=K$
\end{enumerate}

\pyqheader{2026}{16}{MCQ}{2}
The orthogonal trajectories of the family $y=-3x-3+me^{x}$ are given by:
\begin{enumerate}[label=(\Alph*)]
    \item $x=\frac{y}{3}-\frac{1}{9}+ke^{-3y}$
    \item $x=\frac{y}{3}-\frac{1}{9}+ke^{3y}$
    \item $x=-\frac{y}{3}+\frac{1}{9}+ke^{-3y}$
    \item $x=-\frac{y}{3}+\frac{1}{9}+ke^{3y}$
\end{enumerate}

\pyqheader{2026}{18}{MCQ}{2}
Let $\phi:(-\frac{\pi}{2},\frac{\pi}{2})\rightarrow\mathbb{R}$ solve $(\cos x)\frac{dy}{dx}-y=2y^{2}(\sin x-1)\cos x$ with $\phi(0)=\frac{1}{2}$. Then $\phi(\frac{\pi}{4})=$:
\begin{enumerate}[label=(\Alph*)]
    \item 2
    \item $\frac{3+2\sqrt{2}}{2}$
    \item $\frac{1}{2}$
    \item $\frac{1}{\sqrt{2}}$
\end{enumerate}

\pyqheader{2026}{23}{MCQ}{2}
A particular solution of $\frac{d^{2}y}{dx^{2}}+3\frac{dy}{dx}+2y=\sin(e^{x})$ is:
\begin{enumerate}[label=(\Alph*)]
    \item $-e^{-x}\sin(e^{x})$
    \item $-e^{-2x}\sin(e^{x})$
    \item $-e^{-x}(\sin(e^{x})+\cos(e^{x}))$
    \item $-e^{-2x}(\sin(e^{x})+\cos(e^{x}))$
\end{enumerate}

\pyqheader{2026}{35}{MSQ}{2}
Integrating factor(s) for $(2y\cos x-xy\sin x)dx + 2x\cos x\,dy=0$ on $x\in(0,\frac{\pi}{4})$ is/are:
\begin{enumerate}[label=(\Alph*)]
    \item $\frac{1}{xy}$
    \item $xy$
    \item $\sec x$
    \item $\sqrt{\sec x}$
\end{enumerate}

\pyqheader{2026}{45}{NAT}{1}
If $(y^{3}+\alpha xy^{4}-5x+\cos 2y)dx + (3xy^{2}+20x^{2}y^{3}+\beta x\sin 2y)dy=0$ is exact, then $\alpha+\beta=$ \underline{\hspace{1.5cm}}.

\pyqheader{2026}{51}{NAT}{2}
If $y(x)=C_{1}e^{-x}+C_{2}e^{2x}+\alpha xe^{-x}$ is the general solution of $\frac{d^{2}y}{dx^{2}}+\beta\frac{dy}{dx}+\gamma y=-e^{-x}$, then $\alpha(\beta+\gamma)=$ \underline{\hspace{1.5cm}}.

\pyqheader{2025}{2}{MCQ}{1}
For which choice of $N(x,y)$ is $(e^{x}\sin y-2y\sin x)dx + N(x,y)dy=0$ exact?
\begin{enumerate}[label=(\Alph*)]
    \item $e^{x}\sin y+2\cos x$
    \item $e^{x}\cos y+2\cos x$
    \item $e^{x}\cos y+2\sin x$
    \item $e^{x}\sin y+2\sin x$
\end{enumerate}

\pyqheader{2025}{5}{MCQ}{1}
General solution of $\frac{d^{2}y}{dx^{2}}-8\frac{dy}{dx}+16y=2e^{4x}$ is:
\begin{enumerate}[label=(\Alph*)]
    \item $\alpha_{1}e^{4x}+\alpha_{2}xe^{4x}+x^{2}e^{4x}$
    \item $\alpha_{1}e^{4x}+\alpha_{2}xe^{4x}+2x^{2}e^{4x}$
    \item $\alpha_{1}e^{-4x}+\alpha_{2}e^{4x}+2x^{2}e^{4x}$
    \item $\alpha_{1}xe^{-4x}+\alpha_{2}x^{2}e^{-4x}+x^{2}e^{4x}$
\end{enumerate}

\pyqheader{2025}{11}{MCQ}{2}
Orthogonal trajectory to $yx^{2}=\lambda$ passing through $(2,1)$ is:
\begin{enumerate}[label=(\Alph*)]
    \item $y=\frac{x^{2}}{4}$
    \item $x^{2}-2y^{2}=2$
    \item $x-y=1$
    \item $2x-y^{2}=3$
\end{enumerate}

\pyqheader{2025}{12}{MCQ}{2}
Let $\phi$ solve $x\frac{dy}{dx}=(\ln y-\ln x)y$ with $\phi(1)=e^{2}$. Then $\phi(2)=$:
\begin{enumerate}[label=(\Alph*)]
    \item $e^{2}$
    \item $2e^{3}$
    \item $3e^{2}$
    \item $6e^{3}$
\end{enumerate}

\pyqheader{2025}{14}{MCQ}{2}
Let $\phi$ solve $\frac{dy}{dx}=(y-1)(y-3)$ with $\phi(0)=2$. Then:
\begin{enumerate}[label=(\Alph*)]
    \item $\lim_{x\rightarrow\infty}\phi(x)=0$
    \item $\lim_{x\rightarrow\ln\sqrt{2}}\phi(x)=1$
    \item $\lim_{x\rightarrow-\infty}\phi(x)=3$
    \item $\lim_{x\rightarrow\ln(1/\sqrt{2})}\phi(x)=6$
\end{enumerate}

\pyqheader{2025}{29}{MCQ}{2}
Let $\phi(x)=\int_{0}^{x}\phi(t)dt+\sin x$. Then $\lim_{x\rightarrow\pi/2}(\phi(x)-e^{x})=$:
\begin{enumerate}[label=(\Alph*)]
    \item 1
    \item 2
    \item 3
    \item 4
\end{enumerate}

\pyqheader{2025}{32}{MSQ}{2}
If $\mu M dx + \mu N dy = dw$, which of the following is/are TRUE?
\begin{enumerate}[label=(\Alph*)]
    \item $\mu w$ is an integrating factor for $Mdx+Ndy=0$
    \item $\mu w^{2}$ is an integrating factor for $Mdx+Ndy=0$
    \item $w(x,y)=w(0,0)+\int_{0}^{x}(\mu M)(s,0)ds+\int_{0}^{y}(\mu N)(x,t)dt$
    \item $w(x,y)=w(0,0)+\int_{0}^{x}(\mu M)(s,y)ds+\int_{0}^{x}(\mu N)(0,t)dt$
\end{enumerate}

\pyqheader{2025}{33}{MSQ}{2}
Let $\phi:(-1,\infty)\rightarrow(0,\infty)$ solve $\frac{dy}{dx}-2ye^{x}=2e^{x}\sqrt{y}$ with $\phi(0)=1$. Then:
\begin{enumerate}[label=(\Alph*)]
    \item $\phi$ is unbounded
    \item $\lim_{x\rightarrow\ln 2}\phi(x)=(2e-1)^{2}$
    \item $\lim_{x\rightarrow\ln 2}\phi(x)=\sqrt{2e-1}$
    \item $\phi$ is strictly increasing on $(0,\infty)$
\end{enumerate}

\pyqheader{2025}{44}{NAT}{1}
Let $\phi$ solve $4y^{\prime\prime}+16y^{\prime}+25y=0$ with $\phi(0)=1, \phi^{\prime}(0)=-\frac{1}{2}$. Then $\lim_{x\rightarrow\pi/6}e^{2x}\phi(x)=$ \underline{\hspace{1.5cm}} (rounded off to two decimal places).

\pyqheader{2025}{55}{NAT}{2}
Let $\phi$ solve $x^{2}y^{\prime\prime}-xy^{\prime}+y=6x\ln x$ on $(0,\infty)$ with $\phi(1)=-3, \phi(e)=0$. Then $|\phi^{\prime}(1)|=$ \underline{\hspace{1.5cm}} (rounded off to two decimal places).

\pyqheader{2025}{56}{NAT}{2}
Let $\phi$ solve $\frac{dy}{dx}+2xy=2+4x^{2}$ with $\phi(0)=0$. Then $\phi(2)=$ \underline{\hspace{1.5cm}} (rounded off to two decimal places).

\pyqheader{2024}{1}{MCQ}{1}
Let $y_{c}$ solve $\frac{dy}{dx}-y+y^{3}=0, y(0)=c>0$. Then for every $c>0$:
\begin{enumerate}[label=(\Alph*)]
    \item $\lim_{x\rightarrow\infty}y_{c}(x)=0$
    \item $\lim_{x\rightarrow\infty}y_{c}(x)=1$
    \item $\lim_{x\rightarrow\infty}y_{c}(x)=e$
    \item limit does not exist
\end{enumerate}

\pyqheader{2024}{3}{MCQ}{1}
Let $y(x)$ solve $\frac{dy}{dx}=1+y\sec x$ on $(-\pi/2,\pi/2)$ with $y(0)=0$. Then $y(\pi/6)=$:
\begin{enumerate}[label=(\Alph*)]
    \item $\sqrt{3}\log(\frac{3}{2})$
    \item $(\frac{\sqrt{3}}{2})\log(\frac{3}{2})$
    \item $(\frac{\sqrt{3}}{2})\log 3$
    \item $\sqrt{3}\log 3$
\end{enumerate}

\pyqheader{2024}{4}{MCQ}{1}
Differential equation for orthogonal trajectories of $x^{2}+2hxy+y^{2}=1$ is:
\begin{enumerate}[label=(\Alph*)]
    \item $(x^{2}y-y^{3}+y)\frac{dy}{dx}-(xy^{2}-x^{3}+x)=0$
    \item $(x^{2}y-y^{3}+y)\frac{dy}{dx}+(xy^{2}-x^{3}+x)=0$
    \item $(x^{2}y+y^{3}+y)\frac{dy}{dx}-(xy^{2}+x^{3}+x)=0$
    \item $(x^{2}y+y^{3}+y)\frac{dy}{dx}+(xy^{2}+x^{3}+x)=0$
\end{enumerate}

\pyqheader{2024}{10}{MCQ}{1}
Solution of $y^{\prime\prime}+y=g(x)$ with $y(0)=0, y^{\prime}(0)=1$ is:
\begin{enumerate}[label=(\Alph*)]
    \item $y(x)=\sin x-\int_{0}^{x}\sin(x-t)g(t)dt$
    \item $y(x)=\sin x+\int_{0}^{x}\sin(x-t)g(t)dt$
    \item $y(x)=\sin x-\int_{0}^{x}\cos(x-t)g(t)dt$
    \item $y(x)=\sin x+\int_{0}^{x}\cos(x-t)g(t)dt$
\end{enumerate}

\pyqheader{2024}{18}{MCQ}{2}
For $y(8x-9y)dx+2x(x-3y)dy=0$:
\begin{enumerate}[label=(\Alph*)]
    \item Not exact, has $x^{2}$ as integrating factor
    \item Exact and homogeneous
    \item Not exact, does not have $x^2$ as I.F.
    \item Not homogeneous, has $x^2$ as I.F.
\end{enumerate}

\pyqheader{2024}{23}{MCQ}{2}
Let $y_{\alpha}$ solve $\frac{dy}{dx}+2y=\frac{1}{1+x^{2}}$ with $y(0)=\alpha$. Then:
\begin{enumerate}[label=(\Alph*)]
    \item $\lim_{x\rightarrow\infty}y_{\alpha}(x)=0$ for every $\alpha\in\mathbb{R}$
    \item $\lim_{x\rightarrow\infty}y_{\alpha}(x)=1$ for every $\alpha\in\mathbb{R}$
    \item $\exists\alpha$ such that limit exists and $\ne 0,1$
    \item $\exists\alpha$ for which limit does not exist
\end{enumerate}

\pyqheader{2024}{26}{MCQ}{2}
Let $y^{\prime\prime}-2y^{\prime}+y=2e^{x}$. P: $f>0\implies f^{\prime}>0$. Q: $f^{\prime}>0\implies f>0$.
\begin{enumerate}[label=(\Alph*)]
    \item P true, Q false
    \item P false, Q true
    \item Both true
    \item Both false
\end{enumerate}

\pyqheader{2024}{42}{NAT}{1}
Let $y:\mathbb{R}\rightarrow\mathbb{R}$ be the solution to the differential equation $\frac{d^{2}y}{dx^{2}}+2\frac{dy}{dx}+5y=1$ satisfying $y(0)=0$ and $y^{\prime}(0)=1$. Then, $\lim_{x\rightarrow\infty}y(x)$ equals \underline{\hspace{1.5cm}} (rounded off to two decimal places).

\pyqheader{2024}{43}{NAT}{1}
Let $y_{\alpha}$ solve $2y^{\prime\prime}-y^{\prime}-y=0, y(0)=1, y^{\prime}(0)=\alpha>0$. Smallest $\alpha$ for which $y_{\alpha}$ has no critical points in $\mathbb{R}$ is \underline{\hspace{1.5cm}}.

\pyqheader{2024}{51}{NAT}{2}
Number of $\alpha\in(-2\pi,0)\cap\mathbb{Z}$ for which all real solutions of $x^{2}y^{\prime\prime}+\alpha xy^{\prime}+y=0$ approach 0 as $x\rightarrow 0^{+}$ is \underline{\hspace{1.5cm}}.

\pyqheader{2023}{4}{MCQ}{1}
Consider the initial value problem $\frac{dy}{dx}+\alpha y=0, y(0)=1$, where $\alpha\in\mathbb{R}$. Then:
\begin{enumerate}[label=(\Alph*)]
    \item there is an $\alpha$ such that $y(1)=0$
    \item there is a unique $\alpha$ such that $\lim_{x\rightarrow\infty}y(x)=0$
    \item there is NO $\alpha$ such that $y(2)=1$
    \item there is a unique $\alpha$ such that $y(1)=2$
\end{enumerate}

\pyqheader{2023}{22}{MCQ}{2}
Orthogonal trajectory to $x^{2}+y^{2}=2x+4y+k$ passing through $(2,3)$ also passes through:
\begin{enumerate}[label=(\Alph*)]
    \item $(3,4)$
    \item $(-1,1)$
    \item $(1,0)$
    \item $(3,5)$
\end{enumerate}

\pyqheader{2023}{33}{MSQ}{2}
Let $y:(\sqrt{2/3},\infty)\rightarrow\mathbb{R}$ solve $(2x-y)y^{\prime}+(2y-x)=0$ with $y(1)=3$. Then:
\begin{enumerate}[label=(\Alph*)]
    \item $y(3)=1$
    \item $y(2)=4+\sqrt{10}$
    \item $y^{\prime}$ bounded on $(\sqrt{2/3},1)$
    \item $y^{\prime}$ bounded on $(1,\infty)$
\end{enumerate}

\pyqheader{2023}{43}{NAT}{1}
If $y^{\prime\prime}-2y^{\prime}+y=e^{x}, y(0)=0, y^{\prime}(0)=-1/2$, then $y(1)=$ \underline{\hspace{1.5cm}} (rounded off to two decimal places).

\pyqheader{2023}{53}{NAT}{2}
Let $y:(1,\infty)\rightarrow\mathbb{R}$ solve $y^{\prime\prime}-\frac{2y}{(1-x)^{2}}=0$ with $y(2)=1$ and $\lim_{x\rightarrow\infty}y(x)=0$. Then $y(3)=$ \underline{\hspace{1.5cm}} (rounded off to two decimal places).

\pyqheader{2023}{59}{NAT}{2}
Let $S=\{\alpha\in\mathbb{R} : \text{solution of } y^{\prime}=y(2-y), y(0)=\alpha \text{ exists on } [0,\infty)\}$. Then $\min S=$ \underline{\hspace{1.5cm}}.

\pyqheader{2022}{13}{MCQ}{2}
Let $u^{\prime\prime}(x)=\frac{u(x)}{1+x^{2}}$ on $\mathbb{R}$ with $u(0)>0, u^{\prime}(0)>0$. I: $uu^{\prime}$ is increasing on $[0,\infty)$. II: $u$ is increasing on $[0,\infty)$.
\begin{enumerate}[label=(\Alph*)]
    \item Both false
    \item Both true
    \item I false, II true
    \item I true, II false
\end{enumerate}

\pyqheader{2022}{21}{MCQ}{2}
On $(-c,c)$, let $y(x)$ solve $\frac{dy}{dx}=y^{2}-1+\cos x, y(0)=0$. Then:
\begin{enumerate}[label=(\Alph*)]
    \item Local max at origin
    \item Local min at origin
    \item Strictly increasing on $(-\delta,\delta)$
    \item Strictly decreasing on $(-\delta,\delta)$
\end{enumerate}

\pyqheader{2022}{23}{MCQ}{2}
Consider $y^{\prime\prime}+ay^{\prime}+y=\sin x$.
\begin{enumerate}[label=(\Alph*)]
    \item If $a=0$, all solutions unbounded over $\mathbb{R}$
    \item If $a=1$, all solutions unbounded over $(0,\infty)$
    \item If $a=1$, all solutions tend to 0 as $x\rightarrow\infty$
    \item If $a=2$, all solutions bounded over $(-\infty,0)$
\end{enumerate}

\pyqheader{2022}{25}{MCQ}{2}
Let $\phi$ solve $(x^{2}+y^{2})dx-4xy\,dy=0$ with $\phi(1)=1$. Then:
\begin{enumerate}[label=(\Alph*)]
    \item $(3\phi^{2}+x^{2})^{2}=4x$
    \item $(3\phi^{2}-x^{2})^{2}=4x$
    \item $(3\phi^{2}+x^{2})^{2}=4\phi$
    \item $(3\phi^{2}-x^{2})^{2}=4\phi$
\end{enumerate}

\pyqheader{2022}{31}{MSQ}{2}
Let $y^{\prime}=x^{2}+y^{2}+1$ with $y(0)=0$ on $(-c,c)$. Then:
\begin{enumerate}[label=(\Alph*)]
    \item $y(x)$ is odd
    \item $y(x)$ is even
    \item $(y(x))^{2}$ has local min at 0
    \item $(y(x))^{2}$ has local max at 0
\end{enumerate}

\pyqheader{2022}{35}{MSQ}{2}
Consider $y^{\prime\prime}+4y=\sin(ax), x\in\mathbb{R}$.
\begin{enumerate}[label=(\Alph*)]
    \item All solutions periodic for every $a$
    \item All solutions periodic for $a\in\mathbb{R}\setminus\{-2,2\}$
    \item All solutions periodic for $a\in\mathbb{Q}\setminus\{-2,2\}$
    \item If $a\in\mathbb{R}\setminus\mathbb{Q}$, then unique periodic solution
\end{enumerate}

\pyqheader{2022}{43}{NAT}{1}
Let $y^{\prime}+3x^{2}y=x^{2}, y(0)=4$. Then $\lim_{x\rightarrow\infty}y(x)=$ \underline{\hspace{1.5cm}} (rounded off to two decimal places).

\pyqheader{2022}{47}{NAT}{1}
Let $xy^{2}y^{\prime}+y^{3}=\frac{\sin x}{x}, y(\frac{\pi}{2})=0$. Then $y(\frac{5\pi}{2})=$ \underline{\hspace{1.5cm}} (rounded off to two decimal places).

\pyqheader{2022}{49}{NAT}{1}
Let $x^{2}y^{\prime\prime}+xy^{\prime}-4y=x^{2}$ on $(0,\infty)$ such that $\lim_{x\rightarrow 0^{+}}y(x)$ exists and $y(1)=1$. Then $y^{\prime}(1)=$ \underline{\hspace{1.5cm}}.

\pyqheader{2022}{50}{NAT}{1}
Orthogonal trajectory to $y=\frac{c(1-\cos x)}{\sin x}$ passing through $(\frac{\pi}{3},1)$. If $(\frac{\pi}{4},a)$ lies on it, then $a^{4}=$ \underline{\hspace{1.5cm}}.

% ==============================================================================
% SECTION 3: LINEAR ALGEBRA AND ALGEBRA
% ==============================================================================
\section{Linear Algebra and Algebra}

\subsection{Basic Algebra}

\pyqheader{2026}{38}{MSQ}{2}
Let $\binom{n}{r}$ denote the number of ways of choosing $r$ distinct objects out of $n$ distinct objects. Which of the following statements is/are TRUE?
\begin{enumerate}[label=(\Alph*)]
    \item $\sum_{k=0}^{6}(-1)^{k}\binom{13}{2k}=-64$
    \item $\sum_{k=0}^{6}(-1)^{k}\binom{13}{2k+1}=128$
    \item $\sum_{k=0}^{6}(-1)^{k}\binom{13}{2k}=64$
    \item $\sum_{k=0}^{6}(-1)^{k}\binom{13}{2k+1}=-128$
\end{enumerate}

\pyqheader{2026}{58}{NAT}{2}
A fruit shop has 4 different types of bananas. The number of ways in which 12 bananas can be bought with at least one banana from each type, is \underline{\hspace{1.5cm}}.

\pyqheader{2026}{60}{NAT}{2}
$\binom{5}{0}+\binom{6}{1}+\binom{7}{2}+\binom{8}{3}+\binom{9}{4}+\binom{10}{5}+\binom{11}{6}=$ \underline{\hspace{1.5cm}}.

% ------------------------------------------------------------------------------
\subsection{Matrices and Systems of Linear Equations}

\pyqheader{2026}{8}{MCQ}{1}
Let $A$ be an invertible $5\times 5$ real matrix. $B$ is obtained by swapping columns 2 and 3 of $A$, then adding 3 times column 3 to column 5. Then $B^{-1}$ is obtained by:
\begin{enumerate}[label=(\Alph*)]
    \item swapping rows 2 and 3 of $A^{-1}$, then adding 3 times row 3 to row 5
    \item swapping rows 2 and 3 of $A^{-1}$, then adding $-3$ times row 5 to row 3
    \item adding $-3$ times row 3 to row 5, then swapping rows 2 and 3 of $A^{-1}$
    \item adding 3 times row 5 to row 3, then swapping rows 2 and 3 of $A^{-1}$
\end{enumerate}

\pyqheader{2026}{14}{MCQ}{2}
Let $A$ be a $5\times 3$ real matrix of rank 2 and $b$ be a non-zero $5\times 1$ real column vector. Suppose $\begin{pmatrix}1\\2\\3\end{pmatrix}$ and $\begin{pmatrix}4\\5\\6\end{pmatrix}$ are two solutions of the system of linear equations $Ax=b$. Which ONE of the following is also a solution of $Ax=b$?
\begin{enumerate}[label=(\Alph*)]
    \item $\begin{pmatrix}3\\3\\3\end{pmatrix}$
    \item $\begin{pmatrix}5\\7\\9\end{pmatrix}$
    \item $\begin{pmatrix}5\\6\\9\end{pmatrix}$
    \item $\begin{pmatrix}7\\8\\9\end{pmatrix}$
\end{enumerate}

\pyqheader{2026}{25}{MCQ}{2}
Let $P$ be a nonzero $n\times n$ complex matrix such that $v^{*}Pv\ge 0$ for all $v\in\mathbb{C}^{n}$. Then:
\begin{enumerate}[label=(\Alph*)]
    \item All eigenvalues of $P$ are negative
    \item If $Pv_{1}=v_{2}$ and $Pv_{2}=v_{1}$, then $v_{1}=v_{2}$
    \item $v^{*}Pv\ne 0$ for all nonzero $v$
    \item All eigenvalues of $P^{*}$ are negative
\end{enumerate}

\pyqheader{2026}{37}{MSQ}{2}
Let $P=\begin{pmatrix}0&1&0&0&0\\ 1&0&0&0&0\\ 0&0&0&1&0\\ 0&0&0&0&1\\ 0&0&1&0&0\end{pmatrix}\in M_{5}(\mathbb{C})$. Which of the following is/are TRUE?
\begin{enumerate}[label=(\Alph*)]
    \item $\text{nullity}(P-I)\ge 2$
    \item $P$ has 4 distinct eigenvalues in $\mathbb{C}$
    \item $P$ has 3 distinct eigenvalues in $\mathbb{R}$
    \item If $\lambda$ is an eigenvalue, $\exists n\in\mathbb{N}$ such that $\lambda^{n}=1$
\end{enumerate}

\pyqheader{2026}{39}{MSQ}{2}
Let $P\in M_{4,6}(\mathbb{R})$ and $Q\in M_{6,4}(\mathbb{R})$ such that $PQ=0$ and $QP=0$. Then:
\begin{enumerate}[label=(\Alph*)]
    \item $\text{range}(P)\subseteq\text{null}(Q)$ and $\text{range}(Q)\subseteq\text{null}(P)$
    \item $\text{rank}(P)+\text{rank}(Q)\le 4$
    \item If $\text{range}(P)=\text{null}(Q)$, then $\text{rank}(P)+\text{rank}(Q)=4$
    \item $\text{range}(Q)=\text{null}(P)$
\end{enumerate}

\pyqheader{2026}{49}{NAT}{1}
Let $A\in M_{3}(\mathbb{C})$. Suppose $v=\begin{pmatrix}\sqrt{5}i\\ 2i\\ x\end{pmatrix}\in\text{null}(A)\cap\text{range}(A^{T})$. Then $|x|=$ \underline{\hspace{1.5cm}}.

\pyqheader{2026}{50}{NAT}{1}
Let $P$ be a $5\times 5$ real matrix with $\det(P)=2$. Let $Q$ be the matrix of cofactors of $P$. Then $\det(Q)=$ \underline{\hspace{1.5cm}}.

\pyqheader{2026}{56}{NAT}{2}
Let $M\in M_{3}(\mathbb{R})$ satisfy $M\begin{pmatrix}\sqrt{15}\\ \sqrt{15}\\ 0\end{pmatrix}=\begin{pmatrix}\sqrt{3}\\ \sqrt{2}\\ -5\end{pmatrix}$, $M\begin{pmatrix}0\\ 0\\ \sqrt{30}\end{pmatrix}=\begin{pmatrix}\sqrt{15}\\ \sqrt{10}\\ \sqrt{5}\end{pmatrix}$, $M\begin{pmatrix}-\sqrt{15}\\ \sqrt{15}\\ 0\end{pmatrix}=\begin{pmatrix}\sqrt{12}\\ -\sqrt{18}\\ 0\end{pmatrix}$. Then $\det(M)=$ \underline{\hspace{1.5cm}}.

\pyqheader{2026}{57}{NAT}{2}
If $x+2y+2z=1, 2x+3y+z=2, ax+5y+bz=b$ has infinitely many solutions, then $a+b=$ \underline{\hspace{1.5cm}}.

\pyqheader{2026}{59}{NAT}{2}
Let $A\in M_{3}(\mathbb{R})$ represent reflection about $\{(a,b,-a-b): a,b\in\mathbb{R}\}$. Sum of diagonal elements of $A$ is \underline{\hspace{1.5cm}}.

\pyqheader{2025}{7}{MCQ}{1}
Let $M=\begin{pmatrix}2&0&-1\\ 4&1&-4\\ 2&0&x\end{pmatrix}$. If 0 is an eigenvalue of $M$, then $(M^{4}+M)\begin{pmatrix}1\\ 0\\ 1\end{pmatrix}$ equals:
\begin{enumerate}[label=(\Alph*)]
    \item $\begin{pmatrix}1\\ 0\\ 1\end{pmatrix}$
    \item $\begin{pmatrix}2\\ 0\\ 2\end{pmatrix}$
    \item $\begin{pmatrix}5\\ 0\\ 5\end{pmatrix}$
    \item $\begin{pmatrix}17\\ 0\\ 17\end{pmatrix}$
\end{enumerate}

\pyqheader{2025}{15}{MCQ}{2}
Let $M=\begin{pmatrix}6&2&-6&8\\ 5&3&-9&8\\ 3&1&-2&4\end{pmatrix}$ and system $6x_1+2x_2-6x_3+8x_4=8, 5x_1+3x_2-9x_3+8x_4=16, 3x_1+x_2-2x_3+4x_4=32$.
\begin{enumerate}[label=(\Alph*)]
    \item $\text{rank}(M)=3$, system has solution
    \item $\text{rank}(M)=3$, no solution
    \item $\text{rank}(M)=2$, has solution
    \item $\text{rank}(M)=2$, no solution
\end{enumerate}

\pyqheader{2025}{16}{MCQ}{2}
Let $M=\begin{pmatrix}-2&0&0\\ 3&2&3\\ 4&-1&x\end{pmatrix}$ have eigenvalues $-2, 3$. If $M^{3}\begin{pmatrix}0\\ 1\\ 1\end{pmatrix}=\begin{pmatrix}0\\ 125\\ 125\end{pmatrix}$, then:
\begin{enumerate}[label=(\Alph*)]
    \item $x=5$, $M^{2}+M$ invertible
    \item $x\ne 5$, $M^{2}+M$ invertible
    \item $x=5$, $M^{2}+M$ not invertible
    \item $x\ne 5$, $M^{2}+M$ not invertible
\end{enumerate}

\pyqheader{2025}{22}{MCQ}{2}
Let $M=(m_{ij})_{3\times 3}$ be invertible, $\sigma=(1\,2\,3)\in S_3$, and $M_{\sigma}=(m_{i\sigma(j)})$. Then:
\begin{enumerate}[label=(\Alph*)]
    \item $\det(M)=\det(M_\sigma), \text{nullity}(M-M_\sigma)=0$
    \item $\det(M)=-\det(M_\sigma), \text{nullity}(M-M_\sigma)=1$
    \item $\det(M)=\det(M_\sigma), \text{nullity}(M-M_\sigma)=1$
    \item $\det(M)=-\det(M_\sigma), \text{nullity}(M-M_\sigma)=0$
\end{enumerate}

\pyqheader{2024}{9}{MCQ}{1}
P: Unique solution to $Ax=b$ ($m\times n$) $\implies m=n$. Q: $u,v\in V\setminus W \implies u+v\in V\setminus W$.
\begin{enumerate}[label=(\Alph*)]
    \item Both true
    \item P true, Q false
    \item P false, Q true
    \item Both false
\end{enumerate}

\pyqheader{2024}{15}{MCQ}{2}
Let $a=\begin{pmatrix}\frac{1}{\sqrt{3}}\\ -\frac{1}{\sqrt{2}}\\ \frac{1}{\sqrt{6}}\\ 0\end{pmatrix}$. Consider the following two statements.
\begin{itemize}
    \item[P:] The matrix $I_4-aa^{T}$ is invertible.
    \item[Q:] The matrix $I_4-2aa^{T}$ is invertible.
\end{itemize}
Then, which one of the following holds?
\begin{enumerate}[label=(\Alph*)]
    \item P is false but Q is true
    \item P is true but Q is false
    \item Both P and Q are true
    \item Both P and Q are false
\end{enumerate}

\pyqheader{2024}{16}{MCQ}{2}
$A_{6\times 5}, B_{5\times 4}$. P: $\exists Z\ne 0$ s.t. $AZB=0$. Q: $\exists Y\ne 0$ s.t. $BYA=0$.
\begin{enumerate}[label=(\Alph*)]
    \item Both true
    \item P true, Q false
    \item P false, Q true
    \item Both false
\end{enumerate}

\pyqheader{2024}{37}{MSQ}{2}
For $A,B,C\in M_{10}(\mathbb{R})$ with $A=BC$:
\begin{enumerate}[label=(\Alph*)]
    \item $\text{Rowspace}(A)\subseteq\text{Rowspace}(B)$
    \item $\text{Rowspace}(A)\subseteq\text{Rowspace}(C)$
    \item $\text{Colspace}(A)\subseteq\text{Colspace}(B)$
    \item $\text{Colspace}(A)\subseteq\text{Colspace}(C)$
\end{enumerate}

\pyqheader{2024}{44}{NAT}{1}
Let $M=\begin{pmatrix}0&1&2&3\\ 1&0&1&2\\ 2&1&0&1\\ 3&2&1&0\end{pmatrix}$. The $(4,1)$-entry of $M^{-1}$ is \underline{\hspace{1.5cm}} (rounded off to two decimal places).

\pyqheader{2024}{50}{NAT}{1}
Let $M=\begin{pmatrix}0&0&0&0&-1\\ 2&0&0&0&-4\\ 0&2&0&0&0\\ 0&0&2&0&3\\ 0&0&0&2&2\end{pmatrix}$. If $p(x)$ is char. poly, then $p(2)-1=$ \underline{\hspace{1.5cm}}.

\pyqheader{2023}{2}{MCQ}{1}
Let $v_1,\dots,v_9$ be columns of $9\times 9$ matrix $A$ with $\sum a_i v_i = 0$ (not all $a_i=0$). Then $Ax=\sum v_i$ has:
\begin{enumerate}[label=(\Alph*)]
    \item no solution
    \item a unique solution
    \item more than one but only finitely many solutions
    \item infinitely many solutions
\end{enumerate}

\pyqheader{2023}{13}{MCQ}{2}
Let $A=\begin{pmatrix}1&-1&0\\ 0&0&0\\ -2&2&2\end{pmatrix}$ and $B=A^{5}+A^{4}+I_{3}$. Which is NOT an eigenvalue of $B$?
\begin{enumerate}[label=(\Alph*)]
    \item 1
    \item 2
    \item 49
    \item 3
\end{enumerate}

\pyqheader{2023}{14}{MCQ}{2}
The system of linear equations in $x_1, x_2, x_3$:
$\begin{pmatrix}1&1&1\\ 0&-1&1\\ 2&3&\alpha\end{pmatrix}\begin{pmatrix}x_1\\ x_2\\ x_3\end{pmatrix}=\begin{pmatrix}3\\ 1\\ \beta\end{pmatrix}$, where $\alpha, \beta\in\mathbb{R}$, has:
\begin{enumerate}[label=(\Alph*)]
    \item at least one solution for any $\alpha$ and $\beta$
    \item a unique solution for any $\beta$ when $\alpha\ne 1$
    \item NO solution for any $\alpha$ when $\beta\ne 5$
    \item infinitely many solutions for any $\alpha$ when $\beta=5$
\end{enumerate}

\pyqheader{2023}{56}{NAT}{2}
Let $A=\begin{pmatrix}1&1&0&0&1\\ 1&1&1&1&3\\ 1&1&4&4&4\end{pmatrix}$ and $AB=0$ for $5\times 5$ matrix $B$. Max rank of $B$ is \underline{\hspace{1.5cm}}.

\pyqheader{2023}{58}{NAT}{2}
Max number of linearly independent eigenvectors of $\begin{pmatrix}1&1&0&0\\ 2&2&0&0\\ 0&0&3&0\\ 0&0&1&3\end{pmatrix}$ is \underline{\hspace{1.5cm}}.

\pyqheader{2022}{1}{MCQ}{1}
Let $M=\begin{pmatrix}1&1\\ 1&0\end{pmatrix}$. If $M^{8}\begin{pmatrix}1\\ 0\end{pmatrix}=\begin{pmatrix}x\\ y\end{pmatrix}$, then $x=$:
\begin{enumerate}[label=(\Alph*)]
    \item 21
    \item 22
    \item 34
    \item 35
\end{enumerate}

\pyqheader{2022}{2}{MCQ}{1}
Rank of $\begin{pmatrix}1&1&1&0&0&0\\ 1&0&0&1&1&0\\ 0&1&0&1&0&1\\ 0&0&1&0&1&1\end{pmatrix}$ is:
\begin{enumerate}[label=(\Alph*)]
    \item 1
    \item 2
    \item 3
    \item 4
\end{enumerate}

\pyqheader{2022}{16}{MCQ}{2}
Let $P\in M_{4}(\mathbb{R})$ with $P^{4}=0, P^{3}\ne 0$. Which ONE is FALSE?
\begin{enumerate}[label=(\Alph*)]
    \item $\forall v\ne 0, \{v,Pv,P^{2}v,P^{3}v\}$ is linearly independent.
    \item $\text{rank}(P^{k})=4-k$ for $k=1,2,3,4$.
    \item 0 is an eigenvalue of $P$.
    \item If $Q^{4}=0, Q^{3}\ne 0$, then $\exists S$ nonsingular such that $S^{-1}QS=P$.
\end{enumerate}

\pyqheader{2022}{17}{MCQ}{2}
For $(X,Y)=XY-YX$. P: $(X,(Y,Z))+(Y,(Z,X))+(Z,(X,Y))=0$. Q: $(X,(Y,Z))=((X,Y),Z)$.
\begin{enumerate}[label=(\Alph*)]
    \item Both true
    \item P true, Q false
    \item P false, Q true
    \item Both false
\end{enumerate}

\pyqheader{2022}{18}{MCQ}{2}
Consider the system of linear equations:
$x+y+t=4, 2x-4t=7, x+y+z=5, x-3y-z-10t=\lambda$,
where $x,y,z,t$ are variables and $\lambda$ is a constant. Then which one of the following is true?
\begin{enumerate}[label=(\Alph*)]
    \item If $\lambda=1$, then the system has a unique solution.
    \item If $\lambda=2$, then the system has infinitely many solutions.
    \item If $\lambda=1$, then the system has infinitely many solutions.
    \item If $\lambda=2$, then the system has a unique solution.
\end{enumerate}

\pyqheader{2022}{30}{MCQ}{2}
Let $P$ be a $3\times 3$ real matrix having eigenvalues $\lambda_1=0, \lambda_2=1$ and $\lambda_3=-1$. Further, $v_1=\begin{pmatrix}1\\0\\0\end{pmatrix}, v_2=\begin{pmatrix}1\\1\\0\end{pmatrix}$ and $v_3=\begin{pmatrix}1\\0\\1\end{pmatrix}$ are eigenvectors of $P$ corresponding to $\lambda_1, \lambda_2$ and $\lambda_3$, respectively. Then the entry in the first row and the third column of $P$ is:
\begin{enumerate}[label=(\Alph*)]
    \item 0
    \item 1
    \item $-1$
    \item 2
\end{enumerate}

\pyqheader{2022}{53}{NAT}{2}
Let $A=\begin{pmatrix}1&1\\ 0&1\\ -1&1\end{pmatrix}, \|u\|=1, \|v\|=1$. If $Au=\sqrt{2}v$ and $A^{T}v=\sqrt{2}u$, then $|u_1+2\sqrt{2}v_1|=$ \underline{\hspace{1.5cm}}.

\pyqheader{2022}{57}{NAT}{2}
$\det\begin{pmatrix}11&10&10&10\\ 10&11&10&10\\ 10&10&11&10\\ 10&10&10&11\end{pmatrix}=$ \underline{\hspace{1.5cm}}.

% ------------------------------------------------------------------------------
\subsection{Finite Dimensional Vector Spaces}

\pyqheader{2026}{24}{MCQ}{2}
Let $V=\text{span}\{e_1+e_2, e_2+e_3, e_3+e_4+e_5\}\subset\mathbb{R}^{5}$. Let $S=\{P\in M_{2,5}(\mathbb{R}) : \text{null}(P)=V\}$.
\begin{enumerate}[label=(\Alph*)]
    \item $S=\emptyset$
    \item $S=\{0\}$
    \item $S$ is nonempty and $S$ is not a subspace of $M_{2,5}(\mathbb{R})$
    \item $S$ is a nonzero subspace of $M_{2,5}(\mathbb{R})$
\end{enumerate}

\pyqheader{2026}{28}{MCQ}{2}
Let $F:\mathbb{R}^{3}\rightarrow\mathbb{R}^{3}$ be linear preserving distances $d(F(P),F(Q))=d(P,Q)$. Which ONE is FALSE?
\begin{enumerate}[label=(\Alph*)]
    \item $F$ is an injective map.
    \item $F$ is a surjective map.
    \item $\text{vect}(F(P),F(Q))\cdot\text{vect}(F(R),F(S))=\text{vect}(P,Q)\cdot\text{vect}(R,S)$
    \item $\text{vect}(F(P),F(Q))\times\text{vect}(F(R),F(S))=\text{vect}(P,Q)\times\text{vect}(R,S)$
\end{enumerate}

\pyqheader{2026}{40}{MSQ}{2}
Let $V=\{\frac{a+b\sqrt{2}}{c+d\sqrt{2}}: a,b,c,d\in\mathbb{Q}, c^{2}+d^{2}\ne 0\}$. Which of the following is/are FALSE?
\begin{enumerate}[label=(\Alph*)]
    \item $V$ is closed under the usual addition in $\mathbb{R}$.
    \item $V$ is a subspace of the real vector space $\mathbb{R}$.
    \item $V$ is a two dimensional subspace of the vector space $\mathbb{R}$ over $\mathbb{Q}$.
    \item $V$ is a four dimensional subspace of the vector space $\mathbb{R}$ over $\mathbb{Q}$.
\end{enumerate}

\pyqheader{2025}{6}{MCQ}{1}
Define $T:\mathbb{R}^{3}\rightarrow\mathbb{R}^{3}$ by $T(x,y,z)=(x+z, 2x+3y+5z, 2y+2z)$. Then:
\begin{enumerate}[label=(\Alph*)]
    \item $T$ is one-one and $T$ is NOT onto
    \item $T$ is NOT one-one and $T$ is onto
    \item $T$ is one-one and $T$ is onto
    \item $T$ is NOT one-one and $T$ is NOT onto
\end{enumerate}

\pyqheader{2025}{8}{MCQ}{1}
Let $T:P_{2}(\mathbb{R})\rightarrow P_{2}(\mathbb{R})$ be $T(p(x))=p(x+1)$ with matrix $M$ w.r.t. $\{1,x,x^{2}\}$. Then:
\begin{enumerate}[label=(\Alph*)]
    \item The determinant of $M$ is 2
    \item The rank of $M$ is 2
    \item 1 is the only eigenvalue of $M$
    \item The nullity of $M$ is 2
\end{enumerate}

\pyqheader{2025}{24}{MCQ}{2}
Let $f_1,f_2,f_3:\mathbb{R}^{4}\rightarrow\mathbb{R}$ linear with $\ker(f_1)\subseteq\ker(f_2)\cap\ker(f_3)$. Let $T(v)=(f_1(v),f_2(v),f_3(v))$. Nullity of $T$ is:
\begin{enumerate}[label=(\Alph*)]
    \item 1
    \item 2
    \item 3
    \item 4
\end{enumerate}

\pyqheader{2025}{35}{MSQ}{2}
For $u_1=(1,0,0,-1), u_2=(0,2,0,-1), u_3=(0,0,1,-1), u_4=(0,0,0,1)$ in $\mathbb{R}^{4}$:
\begin{enumerate}[label=(\Alph*)]
    \item $\{u_1,u_2,u_3,u_4\}$ is a linearly independent set in $\mathbb{R}^4$
    \item $\{u_1-u_2, u_2-u_3, u_3-u_4, u_4-u_1\}$ is NOT a linearly independent set in $\mathbb{R}^4$
    \item $\{u_1,-u_2, u_3,-u_4\}$ is NOT a linearly independent set in $\mathbb{R}^4$
    \item $\{u_1+u_2, u_2+u_3, u_3+u_4, u_4+u_1\}$ is a linearly independent set in $\mathbb{R}^4$
\end{enumerate}

\pyqheader{2025}{40}{MSQ}{2}
$V_1=\text{span}\{(1,2,3),(1,1,0)\}, V_2=\text{span}\{(1,-1,0)\}, V_3=\text{span}\{(1,1,1)\}, V_4=\text{span}\{(1,3,6)\}, V_5=\text{span}\{(1,0,-3)\}$.
\begin{enumerate}[label=(\Alph*)]
    \item $V_1\cup V_2$ is a subspace of $\mathbb{R}^3$
    \item $V_1\cup V_3$ is a subspace of $\mathbb{R}^3$
    \item $V_1\cup V_4$ is a subspace of $\mathbb{R}^3$
    \item $V_1\cup V_5$ is a subspace of $\mathbb{R}^3$
\end{enumerate}

\pyqheader{2025}{46}{NAT}{1}
$\dim(V_1\cap V_2\cap V_3)$ for $V_1: x+y+2w=0, V_2: 2y+z+w=0, V_3: x+3y+z+3w=0$ in $\mathbb{R}^{4}$ is \underline{\hspace{1.5cm}}.

\pyqheader{2025}{47}{NAT}{1}
Let $T:\mathbb{R}^{3}\rightarrow\mathbb{R}$ linear with $T(1,1,1)=0, T(1,-1,1)=0, T(0,0,1)=16$. Then $T(\frac{1}{2},\frac{2}{3},\frac{3}{4})=$ \underline{\hspace{1.5cm}}.

\pyqheader{2025}{49}{NAT}{1}
$T(p(x))=xp^{\prime}(x)$ and $S(p(x))=(x+1)p^{\prime}(x)$ on $P_{4}(\mathbb{R})$. Nullity of $S\circ T$ is \underline{\hspace{1.5cm}}.

\pyqheader{2025}{58}{NAT}{2}
Dimension of $H=\{p(x)\in P_{3}(\mathbb{R}): xp^{\prime}(x)=3p(x)\}$ is \underline{\hspace{1.5cm}}.

\pyqheader{2024}{13}{MCQ}{2}
Nonzero subspace $V\subseteq M_{7}(\mathbb{C})$ where every non-zero matrix is invertible. $\dim(V)$ over $\mathbb{C}$ is:
\begin{enumerate}[label=(\Alph*)]
    \item 1
    \item 2
    \item 7
    \item 49
\end{enumerate}

\pyqheader{2024}{17}{MCQ}{2}
Let $P_{11}(x)$ be the real vector space of polynomials in $x$ with real coefficients and having degree at most 11, together with the zero polynomial. Let $E=\{s_0(x), s_1(x), \dots, s_{11}(x)\}$ and $F=\{r_0(x), r_1(x), \dots, r_{11}(x)\}$ be subsets of $P_{11}(x)$ having 12 elements each and satisfying $s_0(3)=s_1(3)=\dots=s_{11}(3)=0$ and $r_0(4)=r_1(4)=\dots=r_{11}(4)=1$. Then, which one of the following is TRUE?
\begin{enumerate}[label=(\Alph*)]
    \item Any such $E$ is not necessarily linearly dependent and any such $F$ is not necessarily linearly dependent
    \item Any such $E$ is necessarily linearly dependent but any such $F$ is not necessarily linearly dependent
    \item Any such $E$ is not necessarily linearly dependent but any such $F$ is necessarily linearly dependent
    \item Any such $E$ is necessarily linearly dependent and any such $F$ is necessarily linearly dependent
\end{enumerate}

\pyqheader{2024}{22}{MCQ}{2}
$T(f(x))=f(x)+f^{\prime}(x)$ on $P_{7}(x)$.
\begin{enumerate}[label=(\Alph*)]
    \item $T$ is not a surjective linear transformation
    \item There exists $k\in\mathbb{N}$ such that $T^k$ is the zero linear transformation
    \item 1 and 2 are the eigenvalues of $T$
    \item There exists $r\in\mathbb{N}$ such that $(T-I)^{r}$ is the zero linear transformation
\end{enumerate}

\pyqheader{2024}{45}{NAT}{1}
$\dim(\{f\in P_{12}(x): f(-x)=f(x) \text{ and } f(2024)=0\})=$ \underline{\hspace{1.5cm}}.

\pyqheader{2024}{49}{NAT}{1}
Let $P_7(x)$ be the real vector space of polynomials in $x$ with degree at most 7, together with the zero polynomial. For $r=1,2,\dots,7$, define $s_r(x)=x(x-1)\cdots(x-(r-1))$ and $s_0(x)=1$. If $x^{5}=\sum_{k=0}^{7}\alpha_{5,k}s_k(x)$, where $\alpha_{5,k}\in\mathbb{R}$, then $\alpha_{5,2}$ equals \underline{\hspace{1.5cm}} (rounded off to two decimal places).

\pyqheader{2024}{58}{NAT}{2}
Dimension of continuous piecewise linear functions on $[0,1]$ with nodes $\frac{1}{4}, \frac{1}{2}, \frac{3}{4}$ is \underline{\hspace{1.5cm}}.

\pyqheader{2023}{3}{MCQ}{1}
Which is a subspace of $\mathbb{R}^{3}$?
\begin{enumerate}[label=(\Alph*)]
    \item $\{(x,y,z): (y+z)^{2}+(2x-3y)^{2}=0\}$
    \item $\{(x,y,z): y\in\mathbb{Q}\}$
    \item $\{(x,y,z): yz=0\}$
    \item $\{(x,y,z): x+2y-3z+1=0\}$
\end{enumerate}

\pyqheader{2023}{12}{MCQ}{2}
I: Linear map on $\mathbb{R}^{3}$ with $\text{range}=\text{null}$. II: Linear map on $\mathbb{R}^{2}$ with $\text{range}=\text{null}$.
\begin{enumerate}[label=(\Alph*)]
    \item Both I and II are TRUE
    \item I is TRUE but II is FALSE
    \item II is TRUE but I is FALSE
    \item Both I and II are FALSE
\end{enumerate}

\pyqheader{2023}{15}{MCQ}{2}
I: $\text{span}(S)=\mathbb{R}^{2}\implies\text{span}(S\cap W)=W$. II: $\text{span}(S\cup T)=\text{span}(S)\cup\text{span}(T)$.
\begin{enumerate}[label=(\Alph*)]
    \item Both I and II are TRUE
    \item I is TRUE but II is FALSE
    \item II is TRUE but I is FALSE
    \item Both I and II are FALSE
\end{enumerate}

\pyqheader{2023}{38}{MSQ}{2}
Which of the following is/are true?
\begin{enumerate}[label=(\Alph*)]
    \item Every linear transformation from $\mathbb{R}^2$ to $\mathbb{R}^2$ maps lines onto points or lines
    \item Every surjective linear transformation from $\mathbb{R}^2$ to $\mathbb{R}^2$ maps lines onto lines
    \item Every bijective linear transformation from $\mathbb{R}^2$ to $\mathbb{R}^2$ maps pairs of parallel lines to pairs of parallel lines
    \item Every bijective linear transformation from $\mathbb{R}^2$ to $\mathbb{R}^2$ maps pairs of perpendicular lines to pairs of perpendicular lines
\end{enumerate}

\pyqheader{2023}{39}{MSQ}{2}
Which of the following is/are linear transformations?
\begin{enumerate}[label=(\Alph*)]
    \item $T(x)=\sin x$
    \item $T(A)=\text{trace}(A)$
    \item $T(x,y)=x+y+1$
    \item $T(p(x))=p(1)$
\end{enumerate}

\pyqheader{2023}{42}{NAT}{1}
Rank of $T:P_2(\mathbb{R})\rightarrow P_4(\mathbb{R})$ given by $T(p(x))=p(x^{2})$ is \underline{\hspace{1.5cm}}.

\pyqheader{2023}{57}{NAT}{2}
Dimension of matrices in $M_3(\mathbb{R})$ with each row sum and column sum equal to zero is \underline{\hspace{1.5cm}}.

\pyqheader{2022}{3}{MCQ}{1}
Which is a subspace of $P_{6}(\mathbb{R})$?
\begin{enumerate}[label=(\Alph*)]
    \item $\{f: f(1/2)\notin\mathbb{Q}\}$
    \item $\{f: f(1/2)=1\}$
    \item $\{f: f(1/2)=f(1)\}$
    \item $\{f: f^{\prime}(1/2)=1\}$
\end{enumerate}

\pyqheader{2022}{20}{MCQ}{2}
Let $S=\{P\in M_5(\mathbb{R}): p_{ij}=p_{rs} \text{ for } i+r=j+s\}$. Which is FALSE?
\begin{enumerate}[label=(\Alph*)]
    \item $S$ is a subspace of symmetric matrices
    \item $\dim(S)=5$
    \item $\dim(S)=11$
    \item 5 divides diagonal sum if entries are integers
\end{enumerate}

\pyqheader{2022}{26}{MCQ}{2}
Dimension of $H=\{M\in M_{4}(\mathbb{C}): M^{T}=\overline{M}\}$ over $\mathbb{R}$ is:
\begin{enumerate}[label=(\Alph*)]
    \item 6
    \item 16
    \item 15
    \item 12
\end{enumerate}

\pyqheader{2022}{39}{MSQ}{2}
Let $T:P_{5}\rightarrow\mathbb{R}$ linear with $T(1)=1$ and $T(x(x-1)\cdots(x-k+1))=1$ for $1\le k\le 5$.
\begin{enumerate}[label=(\Alph*)]
    \item $T(x^{4})=15$
    \item $T(x^{3})=5$
    \item $T(x^{4})=14$
    \item $T(x^{3})=3$
\end{enumerate}

\pyqheader{2022}{40}{MSQ}{2}
Let $P$ be a fixed $3\times 3$ matrix with entries in $\mathbb{R}$. Which of the following maps from $M_3(\mathbb{R})$ to $M_3(\mathbb{R})$ is/are linear?
\begin{enumerate}[label=(\Alph*)]
    \item $T_1(M)=MP-PM$
    \item $T_2(M)=M^{2}P-P^{2}M$
    \item $T_3(M)=MP^{2}+P^{2}M$
    \item $T_4(M)=MP^{2}-PM^{2}$
\end{enumerate}

% ------------------------------------------------------------------------------
\subsection{Groups}

\pyqheader{2026}{9}{MCQ}{1}
Which ONE of the following statements is FALSE?
\begin{enumerate}[label=(\Alph*)]
    \item $S_{3}$ is isomorphic to a subgroup of $S_{4}$.
    \item $\mathbb{Z}_{3}$ is isomorphic to a subgroup of $S_{4}$.
    \item $S_{3}$ is isomorphic to a quotient group of $S_{4}$.
    \item $\mathbb{Z}_{6}$ is isomorphic to a quotient group of $S_{4}$.
\end{enumerate}

\pyqheader{2026}{26}{MCQ}{2}
Let $f:G\rightarrow\mathbb{Z}$ be a surjective homomorphism from abelian group $G$, kernel $K$. Which ONE is FALSE?
\begin{enumerate}[label=(\Alph*)]
    \item $\exists\phi:\mathbb{Z}\rightarrow G$ such that $f\circ\phi=\text{id}_{\mathbb{Z}}$
    \item $K$ is isomorphic to a quotient group of $G$
    \item For all non-zero $\alpha:\mathbb{Z}\rightarrow G$ and $\beta:G\rightarrow K$, $\beta\circ\alpha\ne 0$
    \item $G$ has a subgroup isomorphic to $\mathbb{Z}$
\end{enumerate}

\pyqheader{2026}{27}{MCQ}{2}
Which ONE of the following statements is FALSE?
\begin{enumerate}[label=(\Alph*)]
    \item There is a surjective group homomorphism from the additive group of rational numbers to the multiplicative group of all complex roots of unity.
    \item The multiplicative group of complex numbers of modulus one is isomorphic to a quotient group of the additive group of real numbers.
    \item Any group homomorphism from the multiplicative group of nonzero complex numbers into the group of all invertible $2\times 2$ matrices with real entries has nontrivial kernel.
    \item There exists a group homomorphism from the symmetric group on $n$ symbols into the multiplicative group of all invertible $n\times n$ matrices with real entries, which has a trivial kernel.
\end{enumerate}

\pyqheader{2026}{30}{MCQ}{2}
Let $G=\mathcal{P}(\{1,2,3,4\})$ with $A\Delta B=(A\setminus B)\cup(B\setminus A)$. Which ONE is TRUE?
\begin{enumerate}[label=(\Alph*)]
    \item $\{1\}$ is an identity element of $(G,\Delta)$.
    \item $(G,\Delta)$ is an abelian group but not cyclic.
    \item $(G,\Delta)$ is a group and has an element of order 4.
    \item $(G,\Delta)$ is a group and has an element of order 8.
\end{enumerate}

\pyqheader{2026}{47}{NAT}{1}
Let $C_n$ denote a cyclic group having $n$ elements. If there is a surjective group homomorphism from $C_n$ to $C_{30}$, then the total number of such distinct surjective homomorphisms is \underline{\hspace{1.5cm}}.

\pyqheader{2025}{9}{MCQ}{1}
Let $G$ be a finite abelian group of order 10. Let $x_{0}$ be an element of order 2 in $G$. Number of elements in $\{x\in G: x^{3}=x_{0}\}$ is:
\begin{enumerate}[label=(\Alph*)]
    \item 1
    \item 2
    \item 3
    \item 0
\end{enumerate}

\pyqheader{2025}{13}{MCQ}{2}
Let $X=\{x\in S_{4}: x^{3}=\text{id}\}$ and $Y=\{x\in S_{4}: x^{2}\ne\text{id}\}$. If $|X|=m, |Y|=n$:
\begin{enumerate}[label=(\Alph*)]
    \item $m$ is even and $n$ is even
    \item $m$ is odd and $n$ is even
    \item $m$ is even and $n$ is odd
    \item $m$ is odd and $n$ is odd
\end{enumerate}

\pyqheader{2025}{23}{MCQ}{2}
Let $m=$ number of injective homomorphisms from $\mathbb{Z}_3\rightarrow\mathbb{R}/\mathbb{Z}$, and $n=$ number of homomorphisms from $\mathbb{R}/\mathbb{Z}\rightarrow\mathbb{Z}_3$.
\begin{enumerate}[label=(\Alph*)]
    \item $m=2, n=1$
    \item $m=3, n=3$
    \item $m=2, n=3$
    \item $m=1, n=1$
\end{enumerate}

\pyqheader{2025}{39}{MSQ}{2}
For $U(n)=\{x\in\mathbb{Z}_{n}: \gcd(x,n)=1\}$:
\begin{enumerate}[label=(\Alph*)]
    \item $U(8)$ is a cyclic group
    \item $U(5)$ is a cyclic group
    \item $U(12)$ is a cyclic group
    \item $U(9)$ is a cyclic group
\end{enumerate}

\pyqheader{2025}{52}{NAT}{2}
For $\sigma=(1\,2\,3)\in S_{4}$, $|\{\tau\in S_{4}: \tau\sigma\tau^{-1}=\sigma\}|=$ \underline{\hspace{1.5cm}}.

\pyqheader{2025}{59}{NAT}{2}
Let $G$ be an abelian group of order 35. Let $m$ denote the number of elements of order 5 in $G$, and let $n$ denote the number of elements of order 7 in $G$. Then $m+n=$ \underline{\hspace{1.5cm}}.

\pyqheader{2025}{60}{NAT}{2}
The number of surjective group homomorphisms from $S_4$ to $\mathbb{Z}_6$ is equal to \underline{\hspace{1.5cm}}.

\pyqheader{2024}{5}{MCQ}{1}
Let $G$ be a group of order 39 such that it has exactly one subgroup of order 3 and exactly one subgroup of order 13. Then, which one of the following statements is TRUE?
\begin{enumerate}[label=(\Alph*)]
    \item $G$ is necessarily cyclic
    \item $G$ is abelian but need not be cyclic
    \item $G$ need not be abelian
    \item $G$ has 13 elements of order 13
\end{enumerate}

\pyqheader{2024}{6}{MCQ}{1}
Which one of the following is TRUE for $U(n)$ groups?
\begin{enumerate}[label=(\Alph*)]
    \item $U(5)$ is isomorphic to $U(8)$
    \item $U(10)$ is isomorphic to $U(12)$
    \item $U(8)$ is isomorphic to $U(10)$
    \item $U(8)$ is isomorphic to $U(12)$
\end{enumerate}

\pyqheader{2024}{7}{MCQ}{1}
Which one of the following is TRUE for the symmetric group $S_{13}$?
\begin{enumerate}[label=(\Alph*)]
    \item $S_{13}$ has an element of order 42
    \item $S_{13}$ has no element of order 35
    \item $S_{13}$ has an element of order 27
    \item $S_{13}$ has no element of order 60
\end{enumerate}

\pyqheader{2024}{8}{MCQ}{1}
Let $G$ be a finite group containing a non-identity element which is conjugate to its inverse. Then, which one of the following is TRUE?
\begin{enumerate}[label=(\Alph*)]
    \item The order of $G$ is necessarily even
    \item The order of $G$ is not necessarily even
    \item $G$ is necessarily cyclic
    \item $G$ is necessarily abelian but need not be cyclic
\end{enumerate}

\pyqheader{2024}{11}{MCQ}{2}
Which one of the following groups has elements of order 1, 2, 3, 4, 5 but does not have an element of order greater than or equal to 6?
\begin{enumerate}[label=(\Alph*)]
    \item The alternating group $A_6$
    \item The alternating group $A_5$
    \item $S_6$
    \item $S_5$
\end{enumerate}

\pyqheader{2024}{12}{MCQ}{2}
Consider the group $G=\{A\in M_2(\mathbb{R}): AA^{T}=I_2\}$ with respect to matrix multiplication. The center $Z(G)$ has order:
\begin{enumerate}[label=(\Alph*)]
    \item 1
    \item 2
    \item 4
    \item Infinite
\end{enumerate}

\pyqheader{2024}{33}{MSQ}{2}
Which of the following statements is/are TRUE?
\begin{enumerate}[label=(\Alph*)]
    \item The additive group of real numbers is isomorphic to the multiplicative group of positive real numbers
    \item The multiplicative group of nonzero real numbers is isomorphic to the multiplicative group of nonzero complex numbers
    \item The additive group of rational numbers is isomorphic to the multiplicative group of positive rational numbers
    \item The multiplicative group of positive rational numbers is isomorphic to the multiplicative group of nonzero complex numbers
\end{enumerate}

\pyqheader{2024}{36}{MSQ}{2}
Which of the following is/are TRUE for any group $G$?
\begin{enumerate}[label=(\Alph*)]
    \item If $G$ is non-abelian and $Z(G)$ contains more than one element, then the center of the quotient group $G/Z(G)$ contains only one element
    \item If $|G|\ge 2$, then there exists a non-trivial homomorphism from $\mathbb{Z}$ to $G$
    \item If $|G|\ge 2$ and $G$ is non-abelian, then there exists a non-identity isomorphism from $G$ to itself
    \item If $|G|=p^3$, where $p$ is a prime number, then $G$ is necessarily abelian
\end{enumerate}

\pyqheader{2024}{40}{MSQ}{2}
Consider $G=\{m+n\sqrt{2}: m,n\in\mathbb{Z}\}$ as a subgroup of the additive group $\mathbb{R}$. Which of the following statements is/are TRUE?
\begin{enumerate}[label=(\Alph*)]
    \item $G$ is a cyclic subgroup of $\mathbb{R}$ under addition
    \item $G\cap I$ is non-empty for every non-empty open interval $I\subseteq\mathbb{R}$
    \item $G$ is a closed subset of $\mathbb{R}$
    \item $G$ is isomorphic to the group $\mathbb{Z}\times\mathbb{Z}$, where $(m_1,n_1)+(m_2,n_2)=(m_1+m_2, n_1+n_2)$
\end{enumerate}

\pyqheader{2023}{1}{MCQ}{1}
Let $G$ be a finite group. Then $G$ is necessarily a cyclic group if the order of $G$ is:
\begin{enumerate}[label=(\Alph*)]
    \item 4
    \item 7
    \item 6
    \item 10
\end{enumerate}

\pyqheader{2023}{23}{MCQ}{2}
Consider the following statements:
\begin{itemize}
    \item[I.] Every infinite group has infinitely many subgroups.
    \item[II.] There are only finitely many non-isomorphic groups of a given finite order.
\end{itemize}
Then:
\begin{enumerate}[label=(\Alph*)]
    \item both I and II are TRUE
    \item I is TRUE but II is FALSE
    \item I is FALSE but II is TRUE
    \item both I and II are FALSE
\end{enumerate}

\pyqheader{2023}{27}{MCQ}{2}
How many group homomorphisms are there from $\mathbb{Z}_2$ to $S_5$?
\begin{enumerate}[label=(\Alph*)]
    \item 40
    \item 41
    \item 26
    \item 25
\end{enumerate}

\pyqheader{2023}{29}{MCQ}{2}
From the additive group $(\mathbb{Q},+)$, to which one of the following groups does there exist a non-trivial group homomorphism?
\begin{enumerate}[label=(\Alph*)]
    \item $\mathbb{R}^{\times}$, the multiplicative group of non-zero real numbers
    \item $(\mathbb{Z},+)$, the additive group of integers
    \item $\mathbb{Z}_2$, the additive group of integers modulo 2
    \item $\mathbb{Q}^{\times}$, the multiplicative group of non-zero rational numbers
\end{enumerate}

\pyqheader{2023}{32}{MSQ}{2}
Let $0\in A\subseteq\mathbb{Z}$. Which of the following conditions imply that $A$ is a subgroup of $(\mathbb{Z},+)$?
\begin{enumerate}[label=(\Alph*)]
    \item $-2A\subseteq A, A+A=A$
    \item $A=-A, A+2A=A$
    \item $A=-A, A+A=A$
    \item $2A\subseteq A, A+A=A$
\end{enumerate}

\pyqheader{2023}{45}{NAT}{1}
For $\sigma\in S_8$, let $o(\sigma)$ denote the order of $\sigma$. Then $\max\{o(\sigma): \sigma\in S_8\}=$ \underline{\hspace{1.5cm}}.

\pyqheader{2023}{46}{NAT}{1}
The number of group isomorphisms from $\mathbb{Z}_8^{\times}$ onto itself is \underline{\hspace{1.5cm}}.

\pyqheader{2023}{54}{NAT}{2}
The number of permutations in $S_4$ that have exactly 2 cycles in their cycle decomposition is \underline{\hspace{1.5cm}}.

\pyqheader{2022}{4}{MCQ}{1}
Let $G$ be a group of order 2022. Let $H$ and $K$ be subgroups of $G$ of order 337 and 674, respectively. If $H\cup K$ is also a subgroup of $G$, then which one of the following is FALSE?
\begin{enumerate}[label=(\Alph*)]
    \item $H$ is a normal subgroup of $H\cup K$.
    \item The order of $H\cup K$ is 1011.
    \item The order of $H\cup K$ is 674.
    \item $K$ is a normal subgroup of $H\cup K$.
\end{enumerate}

\pyqheader{2022}{19}{MCQ}{2}
Consider the group $(\mathbb{Q},+)$ and its subgroup $(\mathbb{Z},+)$. For the quotient group $\mathbb{Q}/\mathbb{Z}$, which one of the following is FALSE?
\begin{enumerate}[label=(\Alph*)]
    \item $\mathbb{Q}/\mathbb{Z}$ contains a subgroup isomorphic to $(\mathbb{Z},+)$.
    \item There is exactly one group homomorphism from $\mathbb{Q}/\mathbb{Z}$ to $(\mathbb{Q},+)$.
    \item For all $n\in\mathbb{N}$, there exists $g\in\mathbb{Q}/\mathbb{Z}$ such that the order of $g$ is $n$.
    \item $\mathbb{Q}/\mathbb{Z}$ is not a cyclic group.
\end{enumerate}

\pyqheader{2022}{24}{MCQ}{2}
For the group $U_{37}=\{g\in\mathbb{Z}_{37}: \gcd(g,37)=1\}$ under multiplication modulo 37. Which ONE of the following is FALSE?
\begin{enumerate}[label=(\Alph*)]
    \item The set $\{\bar{g}\in U_{37}: \bar{g}=\bar{g}^{-1}\}$ contains exactly 2 elements.
    \item The order of the element 10 in $U_{37}$ is 36.
    \item There is exactly one group homomorphism from $U_{37}$ to $(\mathbb{Z},+)$.
    \item There is exactly one group homomorphism from $U_{37}$ to $(\mathbb{Q},+)$.
\end{enumerate}

\pyqheader{2022}{37}{MSQ}{2}
Let $G$ be a finite group of order at least two. Let $\sigma:G\rightarrow G$ be a bijective group homomorphism such that $\sigma(g)=g\implies g=e$, and $\sigma\circ\sigma=\text{id}$. Which of the following is/are correct?
\begin{enumerate}[label=(\Alph*)]
    \item For each $g\in G$, there exists $h\in G$ such that $h^{-1}\sigma(h)=g$.
    \item There exists $x\in G$ such that $x\sigma(x)\ne e$.
    \item The map $\sigma$ satisfies $\sigma(x)=x^{-1}$ for every $x\in G$.
    \item The order of the group $G$ is an odd number.
\end{enumerate}

\pyqheader{2022}{45}{NAT}{1}
The number of distinct subgroups of $\mathbb{Z}_{999}$ is \underline{\hspace{1.5cm}}.

\pyqheader{2022}{46}{NAT}{1}
The number of elements of order 12 in the symmetric group $S_7$ is equal to \underline{\hspace{1.5cm}}.

\pyqheader{2022}{58}{NAT}{2}
Let $\sigma$ be the permutation in the symmetric group $S_5$ given by $\sigma(1)=2, \sigma(2)=3, \sigma(3)=1, \sigma(4)=5, \sigma(5)=4$. Then the number of elements in the centralizer $N(\sigma)=\{\tau\in S_5: \sigma\circ\tau=\tau\circ\sigma\}$ is equal to \underline{\hspace{1.5cm}}.

% ==============================================================================
% MASTER ANSWER KEY
% ==============================================================================
\newpage
\section{Master Answer Key (2022 -- 2026)}

\begin{center}
\footnotesize
\begin{tabularx}{\textwidth}{ccc|ccc|ccc|ccc|ccc}
\toprule
\multicolumn{3}{c|}{\textbf{JAM 2026}} & \multicolumn{3}{c|}{\textbf{JAM 2025}} & \multicolumn{3}{c|}{\textbf{JAM 2024}} & \multicolumn{3}{c|}{\textbf{JAM 2023}} & \multicolumn{3}{c}{\textbf{JAM 2022}} \\
\textbf{Q\#} & \textbf{Type} & \textbf{Key} & \textbf{Q\#} & \textbf{Type} & \textbf{Key} & \textbf{Q\#} & \textbf{Type} & \textbf{Key} & \textbf{Q\#} & \textbf{Type} & \textbf{Key} & \textbf{Q\#} & \textbf{Type} & \textbf{Key} \\
\midrule
1 & MCQ & D & 1 & MCQ & C & 1 & MCQ & B & 1 & MCQ & B & 1 & MCQ & C \\
2 & MCQ & D & 2 & MCQ & B & 2 & MCQ & A & 2 & MCQ & D & 2 & MCQ & D \\
3 & MCQ & A & 3 & MCQ & A & 3 & MCQ & A & 3 & MCQ & A & 3 & MCQ & C \\
4 & MCQ & D & 4 & MCQ & A & 4 & MCQ & A & 4 & MCQ & D & 4 & MCQ & B \\
5 & MCQ & A & 5 & MCQ & A & 5 & MCQ & A & 5 & MCQ & A & 5 & MCQ & B \\
6 & MCQ & B & 6 & MCQ & D & 6 & MCQ & D & 6 & MCQ & B & 6 & MCQ & D \\
7 & MCQ & B & 7 & MCQ & B & 7 & MCQ & A & 7 & MCQ & A & 7 & MCQ & C \\
8 & MCQ & B & 8 & MCQ & C & 8 & MCQ & A & 8 & MCQ & B & 8 & MCQ & A \\
9 & MCQ & D & 9 & MCQ & A & 9 & MCQ & D & 9 & MCQ & C & 9 & MCQ & A \\
10 & MCQ & MTA & 10 & MCQ & A & 10 & MCQ & B & 10 & MCQ & C & 10 & MCQ & D \\
11 & MCQ & C & 11 & MCQ & B & 11 & MCQ & A & 11 & MCQ & B & 11 & MCQ & C \\
12 & MCQ & C & 12 & MCQ & B & 12 & MCQ & B & 12 & MCQ & C & 12 & MCQ & C \\
13 & MCQ & D & 13 & MCQ & B & 13 & MCQ & A & 13 & MCQ & B & 13 & MCQ & B \\
14 & MCQ & D & 14 & MCQ & C & 14 & MCQ & D & 14 & MCQ & B & 14 & MCQ & D \\
15 & MCQ & D & 15 & MCQ & A & 15 & MCQ & A & 15 & MCQ & D & 15 & MCQ & MTA \\
16 & MCQ & C & 16 & MCQ & B & 16 & MCQ & A & 16 & MCQ & D & 16 & MCQ & A \\
17 & MCQ & D & 17 & MCQ & B & 17 & MCQ & B & 17 & MCQ & B & 17 & MCQ & B \\
18 & MCQ & D & 18 & MCQ & D & 18 & MCQ & A & 18 & MCQ & B & 18 & MCQ & C \\
19 & MCQ & A & 19 & MCQ & D & 19 & MCQ & D & 19 & MCQ & B & 19 & MCQ & A \\
20 & MCQ & A & 20 & MCQ & A & 20 & MCQ & A & 20 & MCQ & C & 20 & MCQ & C \\
21 & MCQ & D & 21 & MCQ & B & 21 & MCQ & A & 21 & MCQ & B & 21 & MCQ & D \\
22 & MCQ & A & 22 & MCQ & C & 22 & MCQ & D & 22 & MCQ & A & 22 & MCQ & C \\
23 & MCQ & B & 23 & MCQ & A & 23 & MCQ & A & 23 & MCQ & A & 23 & MCQ & A \\
24 & MCQ & C & 24 & MCQ & C & 24 & MCQ & A & 24 & MCQ & B & 24 & MCQ & B \\
25 & MCQ & B & 25 & MCQ & A & 25 & MCQ & D & 25 & MCQ & A & 25 & MCQ & B \\
26 & MCQ & C & 26 & MCQ & B & 26 & MCQ & B & 26 & MCQ & D & 26 & MCQ & B \\
27 & MCQ & C & 27 & MCQ & C & 27 & MCQ & A & 27 & MCQ & C & 27 & MCQ & D \\
28 & MCQ & D & 28 & MCQ & B & 28 & MCQ & B & 28 & MCQ & A & 28 & MCQ & B \\
29 & MCQ & A & 29 & MCQ & A & 29 & MCQ & C & 29 & MCQ & A & 29 & MCQ & C \\
30 & MCQ & B & 30 & MCQ & C & 30 & MCQ & A & 30 & MCQ & A & 30 & MCQ & C \\
31 & MSQ & A,C,D & 31 & MSQ & A,C,D & 31 & MSQ & A,B & 31 & MSQ & A,C & 31 & MSQ & A,C \\
32 & MSQ & C,D & 32 & MSQ & A,B,C & 32 & MSQ & A,B & 32 & MSQ & A,B,C & 32 & MSQ & A,B \\
33 & MSQ & A,C & 33 & MSQ & A,B,D & 33 & MSQ & A & 33 & MSQ & B,D & 33 & MSQ & A,B,C \\
34 & MSQ & B & 34 & MSQ & B,D & 34 & MSQ & B,D & 34 & MSQ & A,C,D & 34 & MSQ & A,D \\
35 & MSQ & B,D & 35 & MSQ & A,B & 35 & MSQ & B,C & 35 & MSQ & A,B,D & 35 & MSQ & C,D \\
36 & MSQ & A,C & 36 & MSQ & A,C,D & 36 & MSQ & B,C & 36 & MSQ & A,B,C & 36 & MSQ & C,D \\
37 & MSQ & A,B,D & 37 & MSQ & B,D & 37 & MSQ & B,C & 37 & MSQ & A,B,C & 37 & MSQ & A,C,D \\
38 & MSQ & A & 38 & MSQ & B,C,D & 38 & MSQ & B,D & 38 & MSQ & A,B,C & 38 & MSQ & A,C \\
39 & MSQ & A,B,C & 39 & MSQ & B,D & 39 & MSQ & B,C & 39 & MSQ & B,D & 39 & MSQ & A,B \\
40 & MSQ & B,D & 40 & MSQ & C,D & 40 & MSQ & B,D & 40 & MSQ & A,D & 40 & MSQ & A,C \\
41 & NAT & 4.0 & 41 & NAT & 2.0 & 41 & NAT & 0.25 & 41 & NAT & 0 & 41 & NAT & 2.20 \\
42 & NAT & 1.0 & 42 & NAT & 4.0 & 42 & NAT & 0.20 & 42 & NAT & 3 & 42 & NAT & 14 \\
43 & NAT & 1.0 & 43 & NAT & 2.0 & 43 & NAT & 1.0 & 43 & NAT & 0 & 43 & NAT & 0.33 \\
44 & NAT & 1.0 & 44 & NAT & 1.41 & 44 & NAT & 0.17 & 44 & NAT & 0.50 & 44 & NAT & 0.25 \\
45 & NAT & 8.0 & 45 & NAT & 3.0 & 45 & NAT & 6 & 45 & NAT & 15 & 45 & NAT & 8 \\
46 & NAT & 3.0 & 46 & NAT & 2 & 46 & NAT & 3 & 46 & NAT & 6 & 46 & NAT & 420 \\
47 & NAT & 8 & 47 & NAT & 4.0 & 47 & NAT & 2.0 & 47 & NAT & 0.50 & 47 & NAT & 0 \\
48 & NAT & -4.0 & 48 & NAT & 2.33 & 48 & NAT & 4.56 & 48 & NAT & 4.0 & 48 & NAT & 0.50 \\
49 & NAT & 3.0 & 49 & NAT & 1 & 49 & NAT & 15 & 49 & NAT & 0.50 & 49 & NAT & 2.25 \\
50 & NAT & 16.0 & 50 & NAT & 3.0 & 50 & NAT & 31 & 50 & NAT & 1.0 & 50 & NAT & 2.00 \\
51 & NAT & -1.0 & 51 & NAT & 2.0 & 51 & NAT & 6 & 51 & NAT & 21.0 & 51 & NAT & 18 \\
52 & NAT & 0.25 & 52 & NAT & 3 & 52 & NAT & 0.25 & 52 & NAT & 0.75 & 52 & NAT & 4.16 \\
53 & NAT & 1.25 & 53 & NAT & 9 & 53 & NAT & 2.0 & 53 & NAT & 0.50 & 53 & NAT & 3.00 \\
54 & NAT & 0.45 & 54 & NAT & 5.22 & 54 & NAT & 4.80 & 54 & NAT & 11 & 54 & NAT & 2.00 \\
55 & NAT & 4.0 & 55 & NAT & 1.0 & 55 & NAT & 0.33 & 55 & NAT & 0.50 & 55 & NAT & 3.00 \\
56 & NAT & -1.0 & 56 & NAT & 4.0 & 56 & NAT & 5 & 56 & NAT & 2 & 56 & NAT & 5 \\
57 & NAT & 6.0 & 57 & NAT & MTA & 57 & NAT & 4.5 & 57 & NAT & 4 & 57 & NAT & 41 \\
58 & NAT & 165 & 58 & NAT & 1 & 58 & NAT & 5 & 58 & NAT & 3 & 58 & NAT & 6 \\
59 & NAT & 1.0 & 59 & NAT & 10 & 59 & NAT & 1.0 & 59 & NAT & 0 & 59 & NAT & 2.00 \\
60 & NAT & 924 & 60 & NAT & 0 & 60 & NAT & 1.60 & 60 & NAT & 0.50 & 60 & NAT & 13.40 \\
\bottomrule
\end{tabularx}
\end{center}

\end{document}
